% Include-ready section for 623/supplementary.tex.
% Suggested placement: after "Projective order two: exact enumerator images".
% Usage from the directory containing supplementary.tex:
% % Include-ready section for 623/supplementary.tex.
% Suggested placement: after "Projective order two: exact enumerator images".
% Usage from the directory containing supplementary.tex:
% % Include-ready section for 623/supplementary.tex.
% Suggested placement: after "Projective order two: exact enumerator images".
% Usage from the directory containing supplementary.tex:
% % Include-ready section for 623/supplementary.tex.
% Suggested placement: after "Projective order two: exact enumerator images".
% Usage from the directory containing supplementary.tex:
% \input{orbit2_enumerator_section.tex}
% This file has no document class, preamble, or document environment.
% It uses only the supplement's existing packages, macros, and theorem styles.
% Every new label has the prefix o2gc:.
% Existing abstract/status/count statements elsewhere need separate updating;
% this fragment does not edit or silently replace those statements.

\section{Exact weight-enumerator image of orbit 2}
\label{o2gc:sec}

Orbit~2 is the primitive order-two record
\[
 (q,m,a)=\bigl(2,2,(0,0,1,1,1,1)\bigr).
\]
For the proof, relabel the physical sites in the order
\(B_1,B_2,A_1,A_2,A_3,A_4\). Write
\(\mathsf B=\{B_1,B_2\}\) for the inactive pair and
\(\mathsf A=\{A_1,A_2,A_3,A_4\}\) for the active sites. Its chosen physical
logical involution is
\begin{equation}
 D_{\mathrm{phys}}=D(a;2)=I_{\mathsf B}\otimes D,
 \qquad D=Z_{A_1}Z_{A_2}Z_{A_3}Z_{A_4}.
 \label{o2gc:record}
\end{equation}
Let \(V:\C^2\to(\C^2)^{\otimes6}\) be an isometric encoder, with
code projector \(P=VV^\dagger\), satisfying
\begin{equation}
 D_{\mathrm{phys}}V=VZ_L,\qquad
 V^\dagger EV=\lambda_EI_2\quad(\wt(E)\le2),
 \label{o2gc:KL}
\end{equation}
where \(E\) ranges over Hermitian Pauli strings. For all Pauli weights
put \(\lambda_E=\operatorname{Tr}(PE)/2\); this agrees with the KL scalar
at weights at most two but does not assert scalar compression at higher
weights. We use the supplement's normalization
\[
 A_w=\frac14\sum_{\wt(E)=w}|\operatorname{Tr}(PE)|^2
     =\sum_{\wt(E)=w}\lambda_E^2.
\]

\begin{theorem}[Complete orbit-2 enumerator image]
\label{o2gc:thm-image}
Every exact orbit-2 realization has \(A_3=0\). Its complete
Shor--Laflamme \(A\)-enumerator image is
\begin{equation}
 \boxed{
 \mathcal W_2=
 \left\{\mathcal L(\alpha):0\le\alpha\le1\right\},
 \qquad
 \mathcal L(\alpha)=
 (1,\alpha,1-\alpha,0,11+4\alpha,16-\alpha,3-3\alpha).
 }
 \label{o2gc:image}
\end{equation}
The converse includes every zero-support and singular stratum. No
stabilizer assumption is made.
\end{theorem}

The proof has an analytic reduction and an exact finite certificate.
First, KL and complementary purities force a flat inactive pair and a
specific rank-eight active projector. Mixed KL supplies one common
antiunitary, whose spectral saturation restricts an associated partial
isometry to 32 complex Pauli coordinates. A homogeneous quartic
certificate then forces its weight-three coordinates, and the code's
entire \(A_3\), to vanish. The certificate identity and positivity are
verified with integer and rational arithmetic.

\subsection{Flatness of the inactive pair and the active reflection}
\label{o2gc:subsec-flatness}

Put \(\rho=P/2\). The universal six-qubit rank-two MacWilliams normal
form, with \(x=A_1,y=A_2,z=A_3\), is
\begin{equation}
 (A_0,\ldots,A_6)=
 (1,x,y,z,15-4y-3z,12+3x+4y+3z,4-4x-y-z).
 \label{o2gc:MacWilliams}
\end{equation}
It follows from the quantum MacWilliams transform with factor \(1/32\)
and \(A_j=B_j\) for \(j=0,1,2\). With full six-qubit spin flip
\(\widetilde P=Y^{\otimes6}P^TY^{\otimes6}\), positivity gives
\begin{equation}
 0\le\operatorname{Tr}(P\widetilde P)
 =\frac1{16}\sum_{w=0}^6(-1)^wA_w
 =\frac12(1-A_1-A_2-A_3).
 \label{o2gc:shadow}
\end{equation}

\begin{lemma}[Universal flatness and weight-three support]
\label{o2gc:lem-flatness}
Every encoder in \eqref{o2gc:KL} satisfies
\begin{equation}
 \rho_{\mathsf B}=\frac{I_4}{4},\qquad
 \rho_{\mathsf B A_i}=\frac{I_4}{4}\otimes\rho_{A_i}
 \quad(1\le i\le4),\qquad A_1+A_2+A_3=1.
 \label{o2gc:flatness}
\end{equation}
The only possibly nonzero weight-three coefficients are
\(\lambda_{p_iq_jZ_k}\) on distinct active sites, with
\(p,q\in\{X,Y\}\). Moreover,
\begin{equation}
 \rho_{\mathsf A}=\frac{\Pi}{8},\qquad
 \Pi^\dagger=\Pi=\Pi^2,\qquad \rank\Pi=8,\qquad [\Pi,D]=0,
 \label{o2gc:active-projector}
\end{equation}
and each active parity block of \(\Pi\) has rank four.
\end{lemma}

\begin{proof}
Purify the maximally mixed code with a reference qubit \(R\):
\[
 \ket{\Psi}=
 \frac{\ket0_R\ket{0_L}+\ket1_R\ket{1_L}}{\sqrt2}.
\]
The operator \(G=Z_RI_{\mathsf B}D\) fixes this state. KL says that every
correlation involving a nonidentity reference Pauli and a physical
Pauli of weight at most two vanishes. Multiplication by \(G\) preserves
squared correlations.

Physical coefficients with an odd number of active \(X/Y\) factors
vanish by \(D\)-parity. An active-only coefficient with at least two
\(Z\) factors also vanishes: multiplying by \(G\) gives a
reference-\(Z\) correlation with physical weight \(4-z\le2\), where
\(z\) is the original active \(Z\)-count.
For this purity computation, denote the remaining nonidentity
active-only squared masses by
\[
 \begin{array}{c|cccc}
 \text{mass}&a&u&v&w\\ \hline
 \text{pattern}&Z_i&p_iq_j&p_iq_jZ_k&p_1p_2p_3p_4
 \end{array}
 \qquad(p,q,p_i\in\{X,Y\}).
\]
All displayed sites in a pattern are distinct. Let \(s=a+u+v+w\), and
define the nonnegative masses
\[
 b=\sum_{\substack{F\in\{I,X,Y,Z\}^{\otimes\mathsf B}\\F\ne I}}
       \lambda_F^2,\qquad
 m=\sum_{i=1}^4\sum_{F\ne I}\lambda_{FZ_i}^2.
\]
Here \(F\) in the second sum ranges over all nonidentity spectator
Paulis, including those of spectator weight two.

Purity of a marginal on \(r\) qubits is \(2^{-r}\) times the sum of its
squared Pauli correlations. KL gives
\(\rho_{R\mathsf B}=I_R/2\otimes\rho_{\mathsf B}\).
Equality of complementary purities for \(\mathsf A\) and \(R\mathsf B\)
therefore gives
\begin{equation}
 1+s=2(1+b).
 \label{o2gc:purity-1}
\end{equation}
Let \(h\) be the squared correlation mass on \(R\mathsf A\) with
reference factor \(X\) or \(Y\). Multiplication by \(G\) identifies the
reference-\(Z\) mass with the physical active mass \(1+s\). Comparing
the complementary purities of \(\mathsf B\) and \(R\mathsf A\) gives
\[
 2(1+s)+h=8(1+b),\qquad h=4(1+b).
\]
A surviving reference-\(X/Y\) correlation has odd active transverse
count. If that count is one and its active \(Z\)-count is \(z\), the
term and its \(G\)-multiple have physical weights \(1+z\) and \(4-z\).
At least one weight is at most two, so both correlations vanish.
Thus the surviving transverse count is three. Multiplication by \(G\)
interchanges its physical-weight-three and physical-weight-four
terms, giving
\begin{equation}
 h_3=h_4=\frac h2=2(1+b).
 \label{o2gc:purity-h}
\end{equation}

For each \(i\), compare the complementary marginals
\(\mathsf B A_i\) and \(R(\mathsf A\setminus\{A_i\})\), then sum over
all four active sites. The first marginal's mass is
\[
 1+b+\lambda_{Z_i}^2+\sum_{F\ne I}\lambda_{FZ_i}^2.
\]
On the complementary side, reference identity and reference \(Z\)
together count active strings whose \(i\)th factor is \(I\) or \(Z\).
A string with transverse count \(t\) is therefore counted \(4-t\)
times. Reference-\(X/Y\) physical-weight-three terms appear once and
weight-four terms never appear. Accounting also for the dimensions of
the two marginals gives
\begin{equation}
 2\{4(1+b)+a+m\}
 =4(1+s)-2u-2v-4w+h_3.
 \label{o2gc:purity-2}
\end{equation}
Substitution of \eqref{o2gc:purity-1} and \eqref{o2gc:purity-h} yields
\begin{equation}
 b+m+w=0.
 \label{o2gc:purity-sos}
\end{equation}
Every summand is a sum of real squares. Hence \(b=m=w=0\), every
coefficient they contain vanishes, and \(a+u+v=1\). In particular
\(\rho_{\mathsf B}=I_4/4\). Parity kills spectator-dressed active
\(X/Y\) coefficients, and \(m=0\) kills all spectator-dressed \(Z_i\)
coefficients. This proves the product marginal in
\eqref{o2gc:flatness}.

The full weight-one and weight-two masses are now \(A_1=a\) and
\(A_2=u\). Write \(A_3=v+n\), where \(n\ge0\) collects the remaining
terms supported on two active sites and one spectator site. All other
additional low-weight patterns have vanished. Equation
\eqref{o2gc:shadow} gives \(1+n\le1\), so \(n=0\), \(A_3=v\), and
\(A_1+A_2+A_3=1\).

Finally, KL for all errors on \(\mathsf B\) gives the same marginal
\(I_4/4\) for each logical word. Each word therefore has four equal
Schmidt coefficients \(1/2\). Its active Schmidt space lies in the
corresponding \(D\)-parity sector. Each logical word's active marginal
is its rank-four projector divided by four. Averaging these two
marginals gives \(\Pi/8\), proving
\eqref{o2gc:active-projector}.
\end{proof}

Write the resulting flat parity frames as \(16\times4\) matrices
\(W_0,W_1\):
\begin{equation}
 \ket{z_L}=\frac12\sum_{b=0}^3\ket b_{\mathsf B}\otimes W_z\ket b,
 \qquad W_z^\dagger W_z=I_4,\qquad DW_z=(-1)^zW_z.
 \label{o2gc:frames}
\end{equation}
Their ranges are orthogonal, and
\[
 \Pi=W_0W_0^\dagger+W_1W_1^\dagger
     =4\operatorname{Tr}_{\mathsf B}P.
\]
Set \(K=2\Pi-I_{16}\). Then \(K^\dagger=K\), \(K^2=I_{16}\), and
\([K,D]=0\). For every nonidentity active Pauli \(E\),
\begin{equation}
 \lambda_{I_{\mathsf B}E}
 =\frac18\operatorname{Tr}(\Pi E)
 =\frac1{16}\operatorname{Tr}(KE).
 \label{o2gc:normalization}
\end{equation}
The active support proved above thus gives
\begin{align}
 K&=K_1+K_2+K_3,\label{o2gc:K-support}\\
 K_1&=\sum_{i=1}^4a_iZ_i,\nonumber\\
 K_2&=\sum_{i<j}\sum_{p,q\in\{X,Y\}}u_{ij}^{pq}p_iq_j,\nonumber\\
 K_3&=\sum_{i<j}\sum_{k\notin\{i,j\}}
              \sum_{p,q\in\{X,Y\}}v_{ij;k}^{pq}p_iq_jZ_k.\nonumber
\end{align}
The three sets have \(4,24,48\) real coefficients, respectively. Their
squared masses are
\begin{equation}
 \alpha=\frac1{16}\operatorname{Tr}(K_1^2)=A_1,\qquad
 \beta=\frac1{16}\operatorname{Tr}(K_2^2)=A_2,\qquad
 \nu=\frac1{16}\operatorname{Tr}(K_3^2)=A_3,
 \quad \alpha+\beta+\nu=1.
 \label{o2gc:masses}
\end{equation}
Consequently \eqref{o2gc:MacWilliams} already reduces the enumerator to
\begin{equation}
 (A_0,\ldots,A_6)=
 (1,\alpha,1-\alpha-\nu,\nu,
   11+4\alpha+\nu,16-\alpha-\nu,3-3\alpha).
 \label{o2gc:plane}
\end{equation}
The remaining converse is \(\nu=0\). An arbitrary reflection having
\eqref{o2gc:K-support} need not be completable to a six-qubit code; the
mixed KL constraints are essential in the next step.

\subsection{One common antiunitary from mixed KL}
\label{o2gc:subsec-commonJ}

On the spectator pair choose the fixed basis of
\(\Theta_{\mathsf B}=(Y\otimes Y)\mathcal K\), where \(\mathcal K\)
denotes computational complex conjugation. Explicitly, let \(M_{\mathsf B}\)
have columns
\[
 i\Phi_+,\quad\Phi_-,\quad\Psi_+,\quad i\Psi_-,
 \qquad
 \Phi_\pm=\frac{\ket{00}\pm\ket{11}}{\sqrt2},\quad
 \Psi_\pm=\frac{\ket{01}\pm\ket{10}}{\sqrt2}.
\]
Each column is fixed by \(\Theta_{\mathsf B}\). This is a common
spectator magic basis, not an independently chosen basis for each
observable. In that basis put
\[
 Q_z=W_z\overline{M_{\mathsf B}},\qquad Q=(Q_0\ Q_1).
\]
The conjugate follows from the bipartite expansion
\eqref{o2gc:frames}. We have \(Q^\dagger Q=I_8\) and \(QQ^\dagger=\Pi\).

\begin{lemma}[Common antiunitary equations]
\label{o2gc:lem-commonJ}
For every real single-body active Hamiltonian \(H\), put
\[
 h=Q^\dagger HQ,\qquad
 \mu_H=\frac18\operatorname{Tr}(\Pi H).
\]
Then the same skew-unitary matrix
\begin{equation}
 J_0=\begin{pmatrix}0&I_4\\-I_4&0\end{pmatrix}
 \label{o2gc:J0}
\end{equation}
satisfies
\begin{equation}
 hJ_0+J_0h^T=2\mu_HJ_0.
 \label{o2gc:commonJ}
\end{equation}
\end{lemma}

\begin{proof}
Each single-spectator Pauli \(F\) changes sign under
\(\Theta_{\mathsf B}\). Hence \(M_{\mathsf B}^\dagger F M_{\mathsf B}\)
is imaginary, antisymmetric, and Hermitian. The six such Paulis span
the six-dimensional space of imaginary antisymmetric Hermitian
\(4\times4\) matrices.

For every single-active Pauli \(E\), mixed KL gives, for all six \(F\)
and all \(z,z'\),
\[
 \frac14\operatorname{Tr}\!\left[
 (M_{\mathsf B}^\dagger F M_{\mathsf B})^T
 Q_z^\dagger E Q_{z'}\right]=0.
\]
Orthogonality to this antisymmetric space makes every block
\(Q_z^\dagger EQ_{z'}\) complex symmetric. A diagonal block is also
Hermitian and is therefore real symmetric.

The stronger product-marginal identity in \eqref{o2gc:flatness}, not
merely spectator flatness, gives
\[
 h_{00}+h_{11}=2\mu_HI_4.
\]
Indeed, tracing \((I_{\mathsf B}\otimes H)\rho\) over the active sites
gives the spectator matrix \((h_{00}+h_{11})^T/8\) in this basis;
the product marginal instead gives \(\mu_HI_4/4\).
Thus
\[
 h=\begin{pmatrix}
 R&C\\ \overline C&2\mu_HI_4-R
 \end{pmatrix},
 \qquad R=R^T=\overline R,\qquad C=C^T.
\]
Direct block multiplication proves \eqref{o2gc:commonJ}.
\end{proof}

The antiunitary \(J_0\mathcal K\) squares to \(-I_8\) and reverses
every centered compression \(h-\mu_HI_8\). The same \(J_0\) works for
all twelve single-active Paulis. Separate spectral pairings, or a
condition using only the four \(Z_i\), would not provide the full
hypotheses used in the final certificate.

\subsection{A sparse partial isometry on four active qubits}
\label{o2gc:subsec-T}

Define
\begin{equation}
 \Omega=QJ_0Q^T=Q_0Q_1^T-Q_1Q_0^T,\qquad
 S=Y^{\otimes4},\qquad T=\Omega S.
 \label{o2gc:T-definition}
\end{equation}
Here \(S\) is real, symmetric, and unitary. Using
\(Q^T\overline Q=I_8\) gives
\[
 \Omega^T=-\Omega,\qquad
 \Omega\Omega^\dagger=\Pi,\qquad
 \Omega^\dagger\Omega=\Pi^T.
\]
Therefore
\begin{equation}
 TT^\dagger=\Pi,\qquad
 T^\dagger T=\widetilde\Pi:=S\Pi^TS,\qquad
 ST^TS=-T,\qquad \operatorname{Tr}(TT^\dagger)=8.
 \label{o2gc:partial-isometry}
\end{equation}
The Pauli identity \(SE^TS=(-1)^{\wt(E)}E\) means that only odd
Pauli weights can occur in \(T\).

\begin{lemma}[Exact edge sparsity]
\label{o2gc:lem-edges}
For computational active strings \(x,y\in\{0,1\}^4\),
\[
 d_H(x,y)=1\quad\Longrightarrow\quad\Omega_{xy}=0.
\]
In fact \(\Omega\) is supported only at Hamming distance three, and
\(T\) only at Hamming distance one.
\end{lemma}

\begin{proof}
Parity already makes \(\Omega\) zero between equal-parity strings,
since it pairs \(Q_0\) and \(Q_1\). Fix an even string \(x\), set
\(s_i=(-1)^{x_i}\), \(t=\sum_i a_is_i\), and write \(\bar x\) for
its complement. Consider
\[
 H_x=\sum_i s_iZ_i,\qquad
 r_x=\bra xQ_0,\qquad r_{\bar x}=\bra{\bar x}Q_0.
\]
The allowed reflection support \eqref{o2gc:K-support} gives
\[
 \|r_x\|^2=\Pi_{xx}=\frac{1+t}{2},\qquad
 \|r_{\bar x}\|^2=\frac{1-t}{2},\qquad
 r_xr_{\bar x}^\dagger=\Pi_{x\bar x}=0.
\]
The last equality uses the absence of any term flipping all four
active bits.

The eigenvalue of \(H_x\) on \(y\) is \(4-2d_H(x,y)\). Thus its
even-parity restriction is
\(4(\ket x\bra x-\ket{\bar x}\bra{\bar x})\), and its odd-parity
spectrum is \(\{-2,+2\}\). With \(h_z=Q_z^\dagger H_xQ_z\), this gives
\[
 h_0=4r_x^\dagger r_x-4r_{\bar x}^\dagger r_{\bar x},\qquad
 h_0^Tr_x^T=2(1+t)r_x^T.
\]
Moreover \(\mu_{H_x}=t\). Equation~\eqref{o2gc:commonJ} implies
\[
 h_1=2tI_4-h_0^T,\qquad h_1r_x^T=-2r_x^T.
\]
Set \(w=Q_1r_x^T\). Then
\[
 \bra w(H_x+2I)\ket w=0.
\]
On odd parity \(H_x+2I\) is positive semidefinite, with kernel
consisting exactly of the strings at distance three from \(x\).
Consequently \(w\) has no distance-one components. For even \(x\)
and odd \(y\),
\[
 \Omega_{xy}=r_x(\bra yQ_1)^T=\bra y w,
\]
so the desired entries vanish. If \(r_x=0\), that row already vanishes;
no division is needed. Skew symmetry supplies the reverse ordering.
Only odd distances are possible, so the remaining distance is three.
Right multiplication by \(S\) complements the column string up to a
phase, taking distance three to distance one.
\end{proof}

An operator supported only on distance-one matrix entries has exactly
one \(X/Y\) factor in its Pauli expansion. Combining this with the odd
Pauli weights in \eqref{o2gc:partial-isometry} gives
\begin{align}
 T&=T_1+T_3,\label{o2gc:T-support}\\
 T_1&=\sum_i(t_i^XX_i+t_i^YY_i),\nonumber\\
 T_3&=\sum_i\sum_{\substack{j<k\\j,k\ne i}}
          (t_{i;jk}^XX_i+t_{i;jk}^YY_i)Z_jZ_k.\nonumber
\end{align}
There are \(8+24=32\) complex coefficients.

For each single-active Pauli \(\sigma\), multiply
\eqref{o2gc:commonJ} by \(Q,Q^T\) to obtain
\[
 \Pi\sigma\Omega+\Omega\sigma^T\Pi^T=2\mu_\sigma\Omega.
\]
Since \(S\sigma^TS=-\sigma\), multiplication on the right by \(S\)
yields the physical common-antiunitary equation
\begin{equation}
 \Pi\sigma T-T\sigma\widetilde\Pi=2\mu_\sigma T,\qquad
 \mu_\sigma=\frac18\operatorname{Tr}(\Pi\sigma).
 \label{o2gc:physical-J}
\end{equation}
All steps so far are valid for arbitrary computational zeros and
singular overlap patterns.

\subsection{A homogeneous rigidity theorem}
\label{o2gc:subsec-rigidity}

The remaining finite statement can be formulated independently of
encoders. All operators in this subsection act on the four active
qubits; active identity factors are suppressed. Write
\[
 \tau(M)=\frac1{16}\operatorname{Tr}(M),\qquad
 \mathcal S_1=\{X_i,Y_i:1\le i\le4\},
\]
\[
 \mathcal S_3=\{p_iZ_jZ_k:p\in\{X,Y\},\ j<k,\ i\notin\{j,k\}\}.
\]
For an arbitrary \(T\) in the complex span of
\(\mathcal S_1\cup\mathcal S_3\), let \(T_1,T_3\) denote these two
coordinate parts and put
\begin{equation}
 N=\tau(TT^\dagger),\qquad c_E=\tau(ETT^\dagger).
 \label{o2gc:homogeneous-scalars}
\end{equation}
Each \(c_E\) is real. Define the homogeneous cubic defects
\begin{align}
 F(T)&=TT^\dagger T-2NT,\label{o2gc:defects}\\
 R_\sigma(T)&=TT^\dagger\sigma T-T\sigma T^\dagger T-4c_\sigma T.
 \nonumber
\end{align}
The target and forbidden Pauli sets are
\begin{align}
 \mathcal V_3&=\{p_iq_jZ_k:i<j,\ k\notin\{i,j\},\
                                  p,q\in\{X,Y\}\},\nonumber\\
 \mathcal D_2&=\{Z_iZ_j:i<j\},\nonumber\\
 \mathcal D_3&=\{Z_iZ_jZ_k:i<j<k\},\nonumber\\
 \mathcal F_4&=\{p_iq_jZ_kZ_\ell:i<j,\
                \{k,\ell\}=\{1,2,3,4\}\setminus\{i,j\},\
                p,q\in\{X,Y\}\}.
 \label{o2gc:Pauli-sets}
\end{align}
Their sizes are \(48,6,4,24\), respectively.

\begin{proposition}[Global common-antiunitary rigidity]
\label{o2gc:prop-rigidity}
Suppose \(T\) has the support \eqref{o2gc:T-support} and
\begin{equation}
 F(T)=0,\qquad R_\sigma(T)=0
 \quad\text{for all twelve single-active Paulis},\qquad
 c_E=0\quad(E\in\mathcal D_2\cup\mathcal D_3\cup\mathcal F_4).
 \label{o2gc:rigidity-hypotheses}
\end{equation}
Then \(T_3=0\) and \(c_E=0\) for every \(E\in\mathcal V_3\).
\end{proposition}

For the matrix constructed from an orbit-2 code,
\(N=\tau(\Pi)=1/2\), so \(F(T)=0\) follows from its partial-isometry
identities. Also \(c_\sigma=\mu_\sigma/2\), so
\eqref{o2gc:physical-J} gives \(R_\sigma(T)=0\) with exactly the
factor \(4\) in \eqref{o2gc:defects}. Finally, the support of
\(K=2TT^\dagger-I\) gives all the forbidden zeros in
\eqref{o2gc:rigidity-hypotheses}. Thus the proposition applies once
its certificate is proved.

\subsection{The exact positive quartic certificate}
\label{o2gc:subsec-certificate}

Define the scalar defect contractions
\begin{align}
 g&=\Re\,\tau(T_3^\dagger F(T)),\nonumber\\
 c_E^{13}&=\tau\!\left(E(T_1T_3^\dagger+T_3T_1^\dagger)\right),
 \qquad c_E^{33}=\tau(ET_3T_3^\dagger),\nonumber\\
 f_{\mathcal C}^{13}&=\sum_{E\in\mathcal C}c_Ec_E^{13},\qquad
 f_{\mathcal C}^{33}=\sum_{E\in\mathcal C}c_Ec_E^{33}.
 \label{o2gc:scalar-contractions}
\end{align}
The last two definitions apply to any Pauli set \(\mathcal C\).
All displayed quantities are real; \(g,f_{\mathcal C}^{13}\), and
\(f_{\mathcal C}^{33}\) are homogeneous of degree four.

For a set \(\mathcal I\) of ordered pairs \((\sigma,P')\), where
\(P'\in\mathcal S_1\cup\mathcal S_3\) is a Pauli label, put
\begin{equation}
 J(\mathcal I)=
 \sum_{(\sigma,P')\in\mathcal I}
 \Re\!\left[\overline{t_{P'}}\,
        \tau\!\left(P'\sigma R_\sigma(T)\right)\right].
 \label{o2gc:J-contraction}
\end{equation}
Here \(t_{P'}\) is the coefficient of \(P'\) in \(T\); \(P'\) is not
the code projector \(P\). The five choices needed below are defined
by the following table. In each row, sum over every active site \(i\)
and every indicated Pauli pair; \(P'_i\) is the factor at that site.
\begin{center}
\begin{tabular}{@{}lll@{}}
\toprule
Sum & \(\sigma\) and \(P'\) & Condition at site \(i\)\\
\midrule
\(J_{3,L}\) &
\(\sigma\in\{X_i,Y_i\},\ P'\in\mathcal S_3\) & \(P'_i=Z\)\\
\(J_{3,Z}\) &
\(\sigma=Z_i,\ P'\in\mathcal S_3\) & \(P'_i\in\{X,Y\}\)\\
\(J_{3,\perp}\) &
\(\sigma\in\{X_i,Y_i\},\ P'\in\mathcal S_3\) &
\(P'_i\in\{X,Y\},\ P'_i\ne\sigma_i\)\\
\(J_{1,\perp}\) &
\(\sigma\in\{X_i,Y_i\},\ P'\in\mathcal S_1\) &
\(P'_i\in\{X,Y\},\ P'_i\ne\sigma_i\)\\
\(J_{1,Z}\) &
\(\sigma=Z_i,\ P'\in\mathcal S_1\) & \(P'_i\in\{X,Y\}\)\\
\bottomrule
\end{tabular}
\end{center}
All five sums vanish if the common-antiunitary defects vanish.

The certificate polynomial is
\begin{align}
 \mathcal Q(T)={}&-\sum_{E\in\mathcal V_3}c_E^2
 +\frac5{12}J_{3,L}+\frac56J_{3,Z}
 +\frac7{12}J_{3,\perp}+\frac12g\nonumber\\
 &+\frac32J_{1,\perp}-\frac32J_{1,Z}
 +\frac{17}{6}f_{\mathcal F_4}^{13}
 +\frac{19}{6}f_{\mathcal F_4}^{33}
 +\frac{25}{6}f_{\mathcal D_3}^{13}
 +\frac13f_{\mathcal D_2}^{33}.
 \label{o2gc:Q}
\end{align}

\begin{lemma}[Exact coefficient matrix and positivity]
\label{o2gc:lem-matrix}
Order the 32 Pauli labels in \(\mathcal S_1\cup\mathcal S_3\)
lexicographically using \(I<X<Y<Z\). Let \(z\) contain all
\(528\) monomials \(t_{P'}t_{Q'}\), \(P'\le Q'\), with no
\(\sqrt2\) factors on unequal pairs. There is an integer real
symmetric \(528\times528\) matrix \(M\) such that, identically for
every complex \(T\) in this coordinate space,
\begin{equation}
 \mathcal Q(T)=\frac1{12}z^\dagger Mz.
 \label{o2gc:matrix-identity}
\end{equation}
The \(36\) rows and columns indexed by two \(\mathcal S_1\)
coordinates are zero. On the remaining \(492\) coordinates,
\(M\) has \(111\) blocks and satisfies
\begin{equation}
 M_{\mathrm{remaining}}-4I_{492}>0.
 \label{o2gc:matrix-positive}
\end{equation}
In particular,
\begin{equation}
 \mathcal Q(T)\ge\frac13
 \sum_{\substack{P'\le Q'\\
 P'\in\mathcal S_3\ \text{or}\ Q'\in\mathcal S_3}}
 |t_{P'}t_{Q'}|^2.
 \label{o2gc:coercive}
\end{equation}
\end{lemma}

\begin{proof}[Exact finite verification]
Treat the \(t_{P'}\) and their formal conjugates as independent
variables and expand \(12\mathcal Q\) in bidegree \((2,2)\).
Each monomial has the unique form
\(\overline{t_{P'}t_{Q'}}t_{R'}t_{S'}\), with \(P'\le Q'\) and
\(R'\le S'\). Accumulating all ordered trace contributions into
these unordered pair indices uniquely determines the coefficient
matrix \(M\); it is not a fitted Gram representation.

The exact verifier multiplies Paulis with phases \(1,-1,i,-i\),
expands the normalized traces and real parts in
\eqref{o2gc:scalar-contractions}--\eqref{o2gc:Q}, and checks integral
entries and \(M=M^T\). In particular it proves
\eqref{o2gc:matrix-identity} before imposing any encoder equations.
No floating-point eigenvalues or numerical matrix entries enter
this verification.

Conjugation by a four-qubit Pauli multiplies each \(t_{P'}\) by its
Pauli-conjugation sign and leaves \(\mathcal Q\) invariant. Hence a
matrix entry between the monomials \(t_{P'}t_{Q'}\) and
\(t_{R'}t_{S'}\) can be nonzero only when the Pauli labels of
\(P'Q'\) and \(R'S'\) agree modulo phase. This gives the character
blocks. The verifier checks every off-block zero entry, all \(36\)
zero rows and columns, and every remaining block.

The exact characteristic polynomials of all blocks are summarized
in Table~\ref{o2gc:tab-blocks}. Every linear root is greater than
four. Each quadratic factor \(x^2-bx+c\) has real roots because it
is a factor of the characteristic polynomial of a real symmetric
matrix. Its shifted roots \(y=x-4\) have positive sum \(b-8\) and
positive product \(c-4b+16\), as is checked from every listed factor.
Thus both roots exceed four. This proves
\eqref{o2gc:matrix-positive} once the exact block polynomials have
been verified.

The standard-library verifier also establishes positivity without
factoring characteristic polynomials. For each of all \(111\)
integer blocks, it subtracts \(4I\) and performs exact rational
\(LDL^T\) elimination. Every pivot is strictly positive. Explicitly,
at pivot \(d>0\), the remaining Schur complement is updated by
\[
 A_{ij}\longmapsto A_{ij}-\frac{A_{ik}A_{kj}}{d}.
\]
All entries and pivots are rational, and positive pivots give positive
definiteness by congruence. Every block is checked individually, so
this argument does not assume that one representative block covers
its entire symmetry class.

Finally, the zero coordinates are precisely the
\(\binom{8+1}{2}=36\) products of two weight-one coordinates.
On the remaining coordinates \eqref{o2gc:matrix-positive} implies
\eqref{o2gc:coercive}. The matrix is defined completely by
\eqref{o2gc:Q} and the specified monomial ordering; the accompanying
JSON additionally records every block and every positive pivot.
\end{proof}

\begin{table}[htbp]
\centering
\small
\renewcommand{\arraystretch}{1.2}
\begin{tabular}{@{}lrrp{0.64\linewidth}@{}}
\toprule
Character pattern & Size & Count & Characteristic polynomial of a block of \(M\)\\
\midrule
\(I\)&24&1&
\((x-424)(x-192)^4(x-24)^{11}(x-16)^2
(x^2-48x+192)^3\)\\
\(Z\)&3&4&\((x-384)(x-48)^2\)\\
\(pq\)&3&24&\((x-52)(x^2-460x+47872)\)\\
\(pqZ\)&4&48&\((x-50)(x-46)(x^2-248x+10716)\)\\
\(ZZ\)&8&6&
\((x-32)(x-24)(x^2-408x+9856)(x^2-352x+9216)^2\)\\
\(ZZZ\)&12&4&
\((x-136)(x-48)(x^2-352x+9216)^3(x^2-64x+720)^2\)\\
\(pqZZ\)&4&24&\((x-56)(x-44)(x^2-380x+23936)\)\\
\bottomrule
\end{tabular}
\caption{All exact character blocks after removal of the \(36\) zero
\(\mathcal S_1\mathcal S_1\) coordinates. Here \(p,q\in\{X,Y\}\);
all distinct-site and transverse-axis choices are included.
The counts sum to \(111\), and the sizes with multiplicity sum to
\(492\). Each displayed polynomial is verified for every block in
the corresponding class.}
\label{o2gc:tab-blocks}
\end{table}

\begin{proof}[Proof of Proposition~\ref{o2gc:prop-rigidity}]
Under \eqref{o2gc:rigidity-hypotheses}, \(g\) and all five \(J\)-sums
vanish. Each forbidden-set contraction also vanishes because its
factor \(c_E\) is zero. Hence \eqref{o2gc:Q} and
\eqref{o2gc:coercive} give
\[
 -\sum_{E\in\mathcal V_3}c_E^2
 =\mathcal Q(T)
 \ge\frac13
 \sum_{\substack{P'\le Q'\\
 P'\in\mathcal S_3\ \text{or}\ Q'\in\mathcal S_3}}
 |t_{P'}t_{Q'}|^2\ge0.
\]
The left side is nonpositive, so every \(c_E\),
\(E\in\mathcal V_3\), vanishes. The right side includes
\(|t_{P'}|^4\) for every \(P'\in\mathcal S_3\), so all those
coordinates vanish as well. Thus \(T_3=0\).
\end{proof}

\subsection{Completion of the converse and exact attainment}
\label{o2gc:subsec-finish}

\begin{proof}[Proof of Theorem~\ref{o2gc:thm-image}]
For every exact orbit-2 encoder, the analytic construction gives
\eqref{o2gc:T-support}, the partial-isometry identities, and all twelve
physical equations \eqref{o2gc:physical-J}. As noted after
Proposition~\ref{o2gc:prop-rigidity}, these imply every hypothesis
\eqref{o2gc:rigidity-hypotheses}. The proposition therefore gives
\(c_E=0\) for \(E\in\mathcal V_3\). By
\eqref{o2gc:normalization}, the coefficient of \(E\) in
\(K=2TT^\dagger-I\) is \(2c_E\), so
\begin{equation}
 A_3=\frac1{16}\operatorname{Tr}(K_3^2)
     =4\sum_{E\in\mathcal V_3}c_E^2=0.
 \label{o2gc:A3-zero}
\end{equation}
Equation~\eqref{o2gc:plane} gives \(\mathcal L(\alpha)\), with
\(0\le\alpha\le1\) by \(A_1,A_2\ge0\).

For completeness, the reverse inclusion has the following explicit
split-qubit construction. Start with the standard perfect five-qubit
code with stabilizer generators
\[
 XZZXI,\quad IXZZX,\quad XIXZZ,\quad ZXIXZ.
\]
It is pure: every nonidentity Pauli of weight at most two has zero
compression. Its logical \(Z^{\otimes5}\), multiplied by the first
generator, gives the weight-three logical representative \(-YIIYZ\).
Local Clifford rotations and a permutation therefore give a pure
five-qubit code with a chosen logical \(Z_1Z_2Z_3\).

For \(0\le\theta\le\pi/2\), put \(a=\cos(\theta/2)\),
\(b=\sin(\theta/2)\), and split original site \(1\) by the isometry
\[
 W_\theta\ket0=a\ket{00}+b\ket{11},\qquad
 W_\theta\ket1=a\ket{01}+b\ket{10}.
\]
Its exact identities are
\begin{equation}
 W_\theta W_\theta^\dagger
 =\frac12(I+\cos\theta\,Z\otimes I+\sin\theta\,X\otimes X),
 \qquad (Z\otimes Z)W_\theta=W_\theta Z.
 \label{o2gc:split-construction}
\end{equation}
Take the two split outputs as \(A_1,A_2\), original sites \(2,3\)
as \(A_3,A_4\), and original sites \(4,5\) as \(B_1,B_2\).
The original weight-three logical lifts to \(I_{\mathsf B}D\),
which is precisely \eqref{o2gc:record}.

Compression of any six-qubit error of weight at most two through
\(W_\theta\otimes I^{\otimes4}\) is supported on at most two original
sites. Purity of the five-qubit code therefore proves all KL equations.
If the error has a nonidentity factor on an unsplit site, its scalar
compression is zero. On the split pair its scalar is the identity
coefficient \(\operatorname{Tr}(W_\theta^\dagger E W_\theta)/2\).
Consequently the only nonzero weight-one and weight-two coefficients
are \(\cos\theta\) on \(Z_{A_1}\) and \(\sin\theta\) on
\(X_{A_1}X_{A_2}\). Thus
\[
 A_1=\cos^2\theta=\alpha,\qquad A_2=\sin^2\theta=1-\alpha.
\]
The universal inequality \eqref{o2gc:shadow} gives \(A_3=0\), and
\eqref{o2gc:MacWilliams} gives \(\mathcal L(\alpha)\). Every
\(\alpha\in[0,1]\), including both endpoints, occurs. The quantum
Singleton bound excludes distance at least four at length six and
code dimension two, so these are exact \(((6,2,3))\) realizations.
\end{proof}

\begin{remark}[Scope of the converse]
The proof uses no division by a computational coefficient, row norm,
polarization, Schmidt eigenvalue, or Gram minor. It covers arbitrary
interactions among different \(Z\)-pairs in \(T_3\) and both regular
and singular frame configurations. The necessary common-antiunitary
equations already suffice for this enumerator conclusion; no
sufficiency claim for the matrix completion problem is required.
The theorem determines \(\mathcal W_2\), not all code projectors up to
local unitaries, and does not assert global connectedness or uniqueness
of the split-code family. It also does not prove a stronger
partial-isometry inequality with the common-antiunitary equations
omitted.
\end{remark}

\subsection{Exact certificate files and reproducibility}
\label{o2gc:subsec-reproducibility}

The proof and its exact certificate are distributed together in
\path{orbit2_enumerator_package.zip}. All filenames below are relative
to the root of the unpacked package:
\begin{itemize}
\item \path{orbit2_enumerator_section.tex}:
this include-ready section, using the supplement's existing preamble;
\item \path{orbit2_enumerator_section_preview.tex}:
a standalone document containing the same section and its full preamble;
\item \path{orbit2_exact_certificate.py}:
the standard-library verifier that regenerates the integer matrix
from the quartic polynomial and checks all rational pivots;
\item \path{orbit2_certificate_blocks.json}:
all integer blocks, monomial indices, and positive rational pivots;
\item \path{orbit2_verification_log.txt}:
the successful exact-verification record regenerated for this package;
\item \path{README.txt}:
instructions for compiling, including, and verifying the proof;
\item \path{SHA256SUMS}:
SHA-256 file hashes for both LaTeX sources, the verifier, the certificate
JSON, the verification log, and the README.
\end{itemize}
Run from the unpacked package directory with Python~3, without
third-party packages:
\begin{verbatim}
python3 orbit2_exact_certificate.py
python3 orbit2_exact_certificate.py --json regenerated_blocks.json
\end{verbatim}
The regenerated JSON agrees with the supplied
\path{orbit2_certificate_blocks.json}. The verifier reconstructs the
matrix directly from the explicitly coded polynomial identity; it does
not read the LaTeX files or trust a precomputed positive-semidefinite
matrix. The JSON and log are supporting records rather than runtime
dependencies.
The canonical sparse integer matrix has \(2{,}844\) nonzero entries.
Its SHA-256 under the verifier's specified sparse-entry serialization is
\begin{center}
\small\ttfamily
64b192ccce90edf4b583285307f46265435a51b032f72e6493057138a1820b1d
\end{center}
Discovery of the multipliers can use numerical optimization; the proof
does not. Its finite verification uses exact integer coefficient
expansion and rational \texttt{Fraction} arithmetic. The analytic
encoder reduction is proved above, so this section does not depend
on an unproved converse or on a regular-support assumption from an
earlier problem note.

% This file has no document class, preamble, or document environment.
% It uses only the supplement's existing packages, macros, and theorem styles.
% Every new label has the prefix o2gc:.
% Existing abstract/status/count statements elsewhere need separate updating;
% this fragment does not edit or silently replace those statements.

\section{Exact weight-enumerator image of orbit 2}
\label{o2gc:sec}

Orbit~2 is the primitive order-two record
\[
 (q,m,a)=\bigl(2,2,(0,0,1,1,1,1)\bigr).
\]
For the proof, relabel the physical sites in the order
\(B_1,B_2,A_1,A_2,A_3,A_4\). Write
\(\mathsf B=\{B_1,B_2\}\) for the inactive pair and
\(\mathsf A=\{A_1,A_2,A_3,A_4\}\) for the active sites. Its chosen physical
logical involution is
\begin{equation}
 D_{\mathrm{phys}}=D(a;2)=I_{\mathsf B}\otimes D,
 \qquad D=Z_{A_1}Z_{A_2}Z_{A_3}Z_{A_4}.
 \label{o2gc:record}
\end{equation}
Let \(V:\C^2\to(\C^2)^{\otimes6}\) be an isometric encoder, with
code projector \(P=VV^\dagger\), satisfying
\begin{equation}
 D_{\mathrm{phys}}V=VZ_L,\qquad
 V^\dagger EV=\lambda_EI_2\quad(\wt(E)\le2),
 \label{o2gc:KL}
\end{equation}
where \(E\) ranges over Hermitian Pauli strings. For all Pauli weights
put \(\lambda_E=\operatorname{Tr}(PE)/2\); this agrees with the KL scalar
at weights at most two but does not assert scalar compression at higher
weights. We use the supplement's normalization
\[
 A_w=\frac14\sum_{\wt(E)=w}|\operatorname{Tr}(PE)|^2
     =\sum_{\wt(E)=w}\lambda_E^2.
\]

\begin{theorem}[Complete orbit-2 enumerator image]
\label{o2gc:thm-image}
Every exact orbit-2 realization has \(A_3=0\). Its complete
Shor--Laflamme \(A\)-enumerator image is
\begin{equation}
 \boxed{
 \mathcal W_2=
 \left\{\mathcal L(\alpha):0\le\alpha\le1\right\},
 \qquad
 \mathcal L(\alpha)=
 (1,\alpha,1-\alpha,0,11+4\alpha,16-\alpha,3-3\alpha).
 }
 \label{o2gc:image}
\end{equation}
The converse includes every zero-support and singular stratum. No
stabilizer assumption is made.
\end{theorem}

The proof has an analytic reduction and an exact finite certificate.
First, KL and complementary purities force a flat inactive pair and a
specific rank-eight active projector. Mixed KL supplies one common
antiunitary, whose spectral saturation restricts an associated partial
isometry to 32 complex Pauli coordinates. A homogeneous quartic
certificate then forces its weight-three coordinates, and the code's
entire \(A_3\), to vanish. The certificate identity and positivity are
verified with integer and rational arithmetic.

\subsection{Flatness of the inactive pair and the active reflection}
\label{o2gc:subsec-flatness}

Put \(\rho=P/2\). The universal six-qubit rank-two MacWilliams normal
form, with \(x=A_1,y=A_2,z=A_3\), is
\begin{equation}
 (A_0,\ldots,A_6)=
 (1,x,y,z,15-4y-3z,12+3x+4y+3z,4-4x-y-z).
 \label{o2gc:MacWilliams}
\end{equation}
It follows from the quantum MacWilliams transform with factor \(1/32\)
and \(A_j=B_j\) for \(j=0,1,2\). With full six-qubit spin flip
\(\widetilde P=Y^{\otimes6}P^TY^{\otimes6}\), positivity gives
\begin{equation}
 0\le\operatorname{Tr}(P\widetilde P)
 =\frac1{16}\sum_{w=0}^6(-1)^wA_w
 =\frac12(1-A_1-A_2-A_3).
 \label{o2gc:shadow}
\end{equation}

\begin{lemma}[Universal flatness and weight-three support]
\label{o2gc:lem-flatness}
Every encoder in \eqref{o2gc:KL} satisfies
\begin{equation}
 \rho_{\mathsf B}=\frac{I_4}{4},\qquad
 \rho_{\mathsf B A_i}=\frac{I_4}{4}\otimes\rho_{A_i}
 \quad(1\le i\le4),\qquad A_1+A_2+A_3=1.
 \label{o2gc:flatness}
\end{equation}
The only possibly nonzero weight-three coefficients are
\(\lambda_{p_iq_jZ_k}\) on distinct active sites, with
\(p,q\in\{X,Y\}\). Moreover,
\begin{equation}
 \rho_{\mathsf A}=\frac{\Pi}{8},\qquad
 \Pi^\dagger=\Pi=\Pi^2,\qquad \rank\Pi=8,\qquad [\Pi,D]=0,
 \label{o2gc:active-projector}
\end{equation}
and each active parity block of \(\Pi\) has rank four.
\end{lemma}

\begin{proof}
Purify the maximally mixed code with a reference qubit \(R\):
\[
 \ket{\Psi}=
 \frac{\ket0_R\ket{0_L}+\ket1_R\ket{1_L}}{\sqrt2}.
\]
The operator \(G=Z_RI_{\mathsf B}D\) fixes this state. KL says that every
correlation involving a nonidentity reference Pauli and a physical
Pauli of weight at most two vanishes. Multiplication by \(G\) preserves
squared correlations.

Physical coefficients with an odd number of active \(X/Y\) factors
vanish by \(D\)-parity. An active-only coefficient with at least two
\(Z\) factors also vanishes: multiplying by \(G\) gives a
reference-\(Z\) correlation with physical weight \(4-z\le2\), where
\(z\) is the original active \(Z\)-count.
For this purity computation, denote the remaining nonidentity
active-only squared masses by
\[
 \begin{array}{c|cccc}
 \text{mass}&a&u&v&w\\ \hline
 \text{pattern}&Z_i&p_iq_j&p_iq_jZ_k&p_1p_2p_3p_4
 \end{array}
 \qquad(p,q,p_i\in\{X,Y\}).
\]
All displayed sites in a pattern are distinct. Let \(s=a+u+v+w\), and
define the nonnegative masses
\[
 b=\sum_{\substack{F\in\{I,X,Y,Z\}^{\otimes\mathsf B}\\F\ne I}}
       \lambda_F^2,\qquad
 m=\sum_{i=1}^4\sum_{F\ne I}\lambda_{FZ_i}^2.
\]
Here \(F\) in the second sum ranges over all nonidentity spectator
Paulis, including those of spectator weight two.

Purity of a marginal on \(r\) qubits is \(2^{-r}\) times the sum of its
squared Pauli correlations. KL gives
\(\rho_{R\mathsf B}=I_R/2\otimes\rho_{\mathsf B}\).
Equality of complementary purities for \(\mathsf A\) and \(R\mathsf B\)
therefore gives
\begin{equation}
 1+s=2(1+b).
 \label{o2gc:purity-1}
\end{equation}
Let \(h\) be the squared correlation mass on \(R\mathsf A\) with
reference factor \(X\) or \(Y\). Multiplication by \(G\) identifies the
reference-\(Z\) mass with the physical active mass \(1+s\). Comparing
the complementary purities of \(\mathsf B\) and \(R\mathsf A\) gives
\[
 2(1+s)+h=8(1+b),\qquad h=4(1+b).
\]
A surviving reference-\(X/Y\) correlation has odd active transverse
count. If that count is one and its active \(Z\)-count is \(z\), the
term and its \(G\)-multiple have physical weights \(1+z\) and \(4-z\).
At least one weight is at most two, so both correlations vanish.
Thus the surviving transverse count is three. Multiplication by \(G\)
interchanges its physical-weight-three and physical-weight-four
terms, giving
\begin{equation}
 h_3=h_4=\frac h2=2(1+b).
 \label{o2gc:purity-h}
\end{equation}

For each \(i\), compare the complementary marginals
\(\mathsf B A_i\) and \(R(\mathsf A\setminus\{A_i\})\), then sum over
all four active sites. The first marginal's mass is
\[
 1+b+\lambda_{Z_i}^2+\sum_{F\ne I}\lambda_{FZ_i}^2.
\]
On the complementary side, reference identity and reference \(Z\)
together count active strings whose \(i\)th factor is \(I\) or \(Z\).
A string with transverse count \(t\) is therefore counted \(4-t\)
times. Reference-\(X/Y\) physical-weight-three terms appear once and
weight-four terms never appear. Accounting also for the dimensions of
the two marginals gives
\begin{equation}
 2\{4(1+b)+a+m\}
 =4(1+s)-2u-2v-4w+h_3.
 \label{o2gc:purity-2}
\end{equation}
Substitution of \eqref{o2gc:purity-1} and \eqref{o2gc:purity-h} yields
\begin{equation}
 b+m+w=0.
 \label{o2gc:purity-sos}
\end{equation}
Every summand is a sum of real squares. Hence \(b=m=w=0\), every
coefficient they contain vanishes, and \(a+u+v=1\). In particular
\(\rho_{\mathsf B}=I_4/4\). Parity kills spectator-dressed active
\(X/Y\) coefficients, and \(m=0\) kills all spectator-dressed \(Z_i\)
coefficients. This proves the product marginal in
\eqref{o2gc:flatness}.

The full weight-one and weight-two masses are now \(A_1=a\) and
\(A_2=u\). Write \(A_3=v+n\), where \(n\ge0\) collects the remaining
terms supported on two active sites and one spectator site. All other
additional low-weight patterns have vanished. Equation
\eqref{o2gc:shadow} gives \(1+n\le1\), so \(n=0\), \(A_3=v\), and
\(A_1+A_2+A_3=1\).

Finally, KL for all errors on \(\mathsf B\) gives the same marginal
\(I_4/4\) for each logical word. Each word therefore has four equal
Schmidt coefficients \(1/2\). Its active Schmidt space lies in the
corresponding \(D\)-parity sector. Each logical word's active marginal
is its rank-four projector divided by four. Averaging these two
marginals gives \(\Pi/8\), proving
\eqref{o2gc:active-projector}.
\end{proof}

Write the resulting flat parity frames as \(16\times4\) matrices
\(W_0,W_1\):
\begin{equation}
 \ket{z_L}=\frac12\sum_{b=0}^3\ket b_{\mathsf B}\otimes W_z\ket b,
 \qquad W_z^\dagger W_z=I_4,\qquad DW_z=(-1)^zW_z.
 \label{o2gc:frames}
\end{equation}
Their ranges are orthogonal, and
\[
 \Pi=W_0W_0^\dagger+W_1W_1^\dagger
     =4\operatorname{Tr}_{\mathsf B}P.
\]
Set \(K=2\Pi-I_{16}\). Then \(K^\dagger=K\), \(K^2=I_{16}\), and
\([K,D]=0\). For every nonidentity active Pauli \(E\),
\begin{equation}
 \lambda_{I_{\mathsf B}E}
 =\frac18\operatorname{Tr}(\Pi E)
 =\frac1{16}\operatorname{Tr}(KE).
 \label{o2gc:normalization}
\end{equation}
The active support proved above thus gives
\begin{align}
 K&=K_1+K_2+K_3,\label{o2gc:K-support}\\
 K_1&=\sum_{i=1}^4a_iZ_i,\nonumber\\
 K_2&=\sum_{i<j}\sum_{p,q\in\{X,Y\}}u_{ij}^{pq}p_iq_j,\nonumber\\
 K_3&=\sum_{i<j}\sum_{k\notin\{i,j\}}
              \sum_{p,q\in\{X,Y\}}v_{ij;k}^{pq}p_iq_jZ_k.\nonumber
\end{align}
The three sets have \(4,24,48\) real coefficients, respectively. Their
squared masses are
\begin{equation}
 \alpha=\frac1{16}\operatorname{Tr}(K_1^2)=A_1,\qquad
 \beta=\frac1{16}\operatorname{Tr}(K_2^2)=A_2,\qquad
 \nu=\frac1{16}\operatorname{Tr}(K_3^2)=A_3,
 \quad \alpha+\beta+\nu=1.
 \label{o2gc:masses}
\end{equation}
Consequently \eqref{o2gc:MacWilliams} already reduces the enumerator to
\begin{equation}
 (A_0,\ldots,A_6)=
 (1,\alpha,1-\alpha-\nu,\nu,
   11+4\alpha+\nu,16-\alpha-\nu,3-3\alpha).
 \label{o2gc:plane}
\end{equation}
The remaining converse is \(\nu=0\). An arbitrary reflection having
\eqref{o2gc:K-support} need not be completable to a six-qubit code; the
mixed KL constraints are essential in the next step.

\subsection{One common antiunitary from mixed KL}
\label{o2gc:subsec-commonJ}

On the spectator pair choose the fixed basis of
\(\Theta_{\mathsf B}=(Y\otimes Y)\mathcal K\), where \(\mathcal K\)
denotes computational complex conjugation. Explicitly, let \(M_{\mathsf B}\)
have columns
\[
 i\Phi_+,\quad\Phi_-,\quad\Psi_+,\quad i\Psi_-,
 \qquad
 \Phi_\pm=\frac{\ket{00}\pm\ket{11}}{\sqrt2},\quad
 \Psi_\pm=\frac{\ket{01}\pm\ket{10}}{\sqrt2}.
\]
Each column is fixed by \(\Theta_{\mathsf B}\). This is a common
spectator magic basis, not an independently chosen basis for each
observable. In that basis put
\[
 Q_z=W_z\overline{M_{\mathsf B}},\qquad Q=(Q_0\ Q_1).
\]
The conjugate follows from the bipartite expansion
\eqref{o2gc:frames}. We have \(Q^\dagger Q=I_8\) and \(QQ^\dagger=\Pi\).

\begin{lemma}[Common antiunitary equations]
\label{o2gc:lem-commonJ}
For every real single-body active Hamiltonian \(H\), put
\[
 h=Q^\dagger HQ,\qquad
 \mu_H=\frac18\operatorname{Tr}(\Pi H).
\]
Then the same skew-unitary matrix
\begin{equation}
 J_0=\begin{pmatrix}0&I_4\\-I_4&0\end{pmatrix}
 \label{o2gc:J0}
\end{equation}
satisfies
\begin{equation}
 hJ_0+J_0h^T=2\mu_HJ_0.
 \label{o2gc:commonJ}
\end{equation}
\end{lemma}

\begin{proof}
Each single-spectator Pauli \(F\) changes sign under
\(\Theta_{\mathsf B}\). Hence \(M_{\mathsf B}^\dagger F M_{\mathsf B}\)
is imaginary, antisymmetric, and Hermitian. The six such Paulis span
the six-dimensional space of imaginary antisymmetric Hermitian
\(4\times4\) matrices.

For every single-active Pauli \(E\), mixed KL gives, for all six \(F\)
and all \(z,z'\),
\[
 \frac14\operatorname{Tr}\!\left[
 (M_{\mathsf B}^\dagger F M_{\mathsf B})^T
 Q_z^\dagger E Q_{z'}\right]=0.
\]
Orthogonality to this antisymmetric space makes every block
\(Q_z^\dagger EQ_{z'}\) complex symmetric. A diagonal block is also
Hermitian and is therefore real symmetric.

The stronger product-marginal identity in \eqref{o2gc:flatness}, not
merely spectator flatness, gives
\[
 h_{00}+h_{11}=2\mu_HI_4.
\]
Indeed, tracing \((I_{\mathsf B}\otimes H)\rho\) over the active sites
gives the spectator matrix \((h_{00}+h_{11})^T/8\) in this basis;
the product marginal instead gives \(\mu_HI_4/4\).
Thus
\[
 h=\begin{pmatrix}
 R&C\\ \overline C&2\mu_HI_4-R
 \end{pmatrix},
 \qquad R=R^T=\overline R,\qquad C=C^T.
\]
Direct block multiplication proves \eqref{o2gc:commonJ}.
\end{proof}

The antiunitary \(J_0\mathcal K\) squares to \(-I_8\) and reverses
every centered compression \(h-\mu_HI_8\). The same \(J_0\) works for
all twelve single-active Paulis. Separate spectral pairings, or a
condition using only the four \(Z_i\), would not provide the full
hypotheses used in the final certificate.

\subsection{A sparse partial isometry on four active qubits}
\label{o2gc:subsec-T}

Define
\begin{equation}
 \Omega=QJ_0Q^T=Q_0Q_1^T-Q_1Q_0^T,\qquad
 S=Y^{\otimes4},\qquad T=\Omega S.
 \label{o2gc:T-definition}
\end{equation}
Here \(S\) is real, symmetric, and unitary. Using
\(Q^T\overline Q=I_8\) gives
\[
 \Omega^T=-\Omega,\qquad
 \Omega\Omega^\dagger=\Pi,\qquad
 \Omega^\dagger\Omega=\Pi^T.
\]
Therefore
\begin{equation}
 TT^\dagger=\Pi,\qquad
 T^\dagger T=\widetilde\Pi:=S\Pi^TS,\qquad
 ST^TS=-T,\qquad \operatorname{Tr}(TT^\dagger)=8.
 \label{o2gc:partial-isometry}
\end{equation}
The Pauli identity \(SE^TS=(-1)^{\wt(E)}E\) means that only odd
Pauli weights can occur in \(T\).

\begin{lemma}[Exact edge sparsity]
\label{o2gc:lem-edges}
For computational active strings \(x,y\in\{0,1\}^4\),
\[
 d_H(x,y)=1\quad\Longrightarrow\quad\Omega_{xy}=0.
\]
In fact \(\Omega\) is supported only at Hamming distance three, and
\(T\) only at Hamming distance one.
\end{lemma}

\begin{proof}
Parity already makes \(\Omega\) zero between equal-parity strings,
since it pairs \(Q_0\) and \(Q_1\). Fix an even string \(x\), set
\(s_i=(-1)^{x_i}\), \(t=\sum_i a_is_i\), and write \(\bar x\) for
its complement. Consider
\[
 H_x=\sum_i s_iZ_i,\qquad
 r_x=\bra xQ_0,\qquad r_{\bar x}=\bra{\bar x}Q_0.
\]
The allowed reflection support \eqref{o2gc:K-support} gives
\[
 \|r_x\|^2=\Pi_{xx}=\frac{1+t}{2},\qquad
 \|r_{\bar x}\|^2=\frac{1-t}{2},\qquad
 r_xr_{\bar x}^\dagger=\Pi_{x\bar x}=0.
\]
The last equality uses the absence of any term flipping all four
active bits.

The eigenvalue of \(H_x\) on \(y\) is \(4-2d_H(x,y)\). Thus its
even-parity restriction is
\(4(\ket x\bra x-\ket{\bar x}\bra{\bar x})\), and its odd-parity
spectrum is \(\{-2,+2\}\). With \(h_z=Q_z^\dagger H_xQ_z\), this gives
\[
 h_0=4r_x^\dagger r_x-4r_{\bar x}^\dagger r_{\bar x},\qquad
 h_0^Tr_x^T=2(1+t)r_x^T.
\]
Moreover \(\mu_{H_x}=t\). Equation~\eqref{o2gc:commonJ} implies
\[
 h_1=2tI_4-h_0^T,\qquad h_1r_x^T=-2r_x^T.
\]
Set \(w=Q_1r_x^T\). Then
\[
 \bra w(H_x+2I)\ket w=0.
\]
On odd parity \(H_x+2I\) is positive semidefinite, with kernel
consisting exactly of the strings at distance three from \(x\).
Consequently \(w\) has no distance-one components. For even \(x\)
and odd \(y\),
\[
 \Omega_{xy}=r_x(\bra yQ_1)^T=\bra y w,
\]
so the desired entries vanish. If \(r_x=0\), that row already vanishes;
no division is needed. Skew symmetry supplies the reverse ordering.
Only odd distances are possible, so the remaining distance is three.
Right multiplication by \(S\) complements the column string up to a
phase, taking distance three to distance one.
\end{proof}

An operator supported only on distance-one matrix entries has exactly
one \(X/Y\) factor in its Pauli expansion. Combining this with the odd
Pauli weights in \eqref{o2gc:partial-isometry} gives
\begin{align}
 T&=T_1+T_3,\label{o2gc:T-support}\\
 T_1&=\sum_i(t_i^XX_i+t_i^YY_i),\nonumber\\
 T_3&=\sum_i\sum_{\substack{j<k\\j,k\ne i}}
          (t_{i;jk}^XX_i+t_{i;jk}^YY_i)Z_jZ_k.\nonumber
\end{align}
There are \(8+24=32\) complex coefficients.

For each single-active Pauli \(\sigma\), multiply
\eqref{o2gc:commonJ} by \(Q,Q^T\) to obtain
\[
 \Pi\sigma\Omega+\Omega\sigma^T\Pi^T=2\mu_\sigma\Omega.
\]
Since \(S\sigma^TS=-\sigma\), multiplication on the right by \(S\)
yields the physical common-antiunitary equation
\begin{equation}
 \Pi\sigma T-T\sigma\widetilde\Pi=2\mu_\sigma T,\qquad
 \mu_\sigma=\frac18\operatorname{Tr}(\Pi\sigma).
 \label{o2gc:physical-J}
\end{equation}
All steps so far are valid for arbitrary computational zeros and
singular overlap patterns.

\subsection{A homogeneous rigidity theorem}
\label{o2gc:subsec-rigidity}

The remaining finite statement can be formulated independently of
encoders. All operators in this subsection act on the four active
qubits; active identity factors are suppressed. Write
\[
 \tau(M)=\frac1{16}\operatorname{Tr}(M),\qquad
 \mathcal S_1=\{X_i,Y_i:1\le i\le4\},
\]
\[
 \mathcal S_3=\{p_iZ_jZ_k:p\in\{X,Y\},\ j<k,\ i\notin\{j,k\}\}.
\]
For an arbitrary \(T\) in the complex span of
\(\mathcal S_1\cup\mathcal S_3\), let \(T_1,T_3\) denote these two
coordinate parts and put
\begin{equation}
 N=\tau(TT^\dagger),\qquad c_E=\tau(ETT^\dagger).
 \label{o2gc:homogeneous-scalars}
\end{equation}
Each \(c_E\) is real. Define the homogeneous cubic defects
\begin{align}
 F(T)&=TT^\dagger T-2NT,\label{o2gc:defects}\\
 R_\sigma(T)&=TT^\dagger\sigma T-T\sigma T^\dagger T-4c_\sigma T.
 \nonumber
\end{align}
The target and forbidden Pauli sets are
\begin{align}
 \mathcal V_3&=\{p_iq_jZ_k:i<j,\ k\notin\{i,j\},\
                                  p,q\in\{X,Y\}\},\nonumber\\
 \mathcal D_2&=\{Z_iZ_j:i<j\},\nonumber\\
 \mathcal D_3&=\{Z_iZ_jZ_k:i<j<k\},\nonumber\\
 \mathcal F_4&=\{p_iq_jZ_kZ_\ell:i<j,\
                \{k,\ell\}=\{1,2,3,4\}\setminus\{i,j\},\
                p,q\in\{X,Y\}\}.
 \label{o2gc:Pauli-sets}
\end{align}
Their sizes are \(48,6,4,24\), respectively.

\begin{proposition}[Global common-antiunitary rigidity]
\label{o2gc:prop-rigidity}
Suppose \(T\) has the support \eqref{o2gc:T-support} and
\begin{equation}
 F(T)=0,\qquad R_\sigma(T)=0
 \quad\text{for all twelve single-active Paulis},\qquad
 c_E=0\quad(E\in\mathcal D_2\cup\mathcal D_3\cup\mathcal F_4).
 \label{o2gc:rigidity-hypotheses}
\end{equation}
Then \(T_3=0\) and \(c_E=0\) for every \(E\in\mathcal V_3\).
\end{proposition}

For the matrix constructed from an orbit-2 code,
\(N=\tau(\Pi)=1/2\), so \(F(T)=0\) follows from its partial-isometry
identities. Also \(c_\sigma=\mu_\sigma/2\), so
\eqref{o2gc:physical-J} gives \(R_\sigma(T)=0\) with exactly the
factor \(4\) in \eqref{o2gc:defects}. Finally, the support of
\(K=2TT^\dagger-I\) gives all the forbidden zeros in
\eqref{o2gc:rigidity-hypotheses}. Thus the proposition applies once
its certificate is proved.

\subsection{The exact positive quartic certificate}
\label{o2gc:subsec-certificate}

Define the scalar defect contractions
\begin{align}
 g&=\Re\,\tau(T_3^\dagger F(T)),\nonumber\\
 c_E^{13}&=\tau\!\left(E(T_1T_3^\dagger+T_3T_1^\dagger)\right),
 \qquad c_E^{33}=\tau(ET_3T_3^\dagger),\nonumber\\
 f_{\mathcal C}^{13}&=\sum_{E\in\mathcal C}c_Ec_E^{13},\qquad
 f_{\mathcal C}^{33}=\sum_{E\in\mathcal C}c_Ec_E^{33}.
 \label{o2gc:scalar-contractions}
\end{align}
The last two definitions apply to any Pauli set \(\mathcal C\).
All displayed quantities are real; \(g,f_{\mathcal C}^{13}\), and
\(f_{\mathcal C}^{33}\) are homogeneous of degree four.

For a set \(\mathcal I\) of ordered pairs \((\sigma,P')\), where
\(P'\in\mathcal S_1\cup\mathcal S_3\) is a Pauli label, put
\begin{equation}
 J(\mathcal I)=
 \sum_{(\sigma,P')\in\mathcal I}
 \Re\!\left[\overline{t_{P'}}\,
        \tau\!\left(P'\sigma R_\sigma(T)\right)\right].
 \label{o2gc:J-contraction}
\end{equation}
Here \(t_{P'}\) is the coefficient of \(P'\) in \(T\); \(P'\) is not
the code projector \(P\). The five choices needed below are defined
by the following table. In each row, sum over every active site \(i\)
and every indicated Pauli pair; \(P'_i\) is the factor at that site.
\begin{center}
\begin{tabular}{@{}lll@{}}
\toprule
Sum & \(\sigma\) and \(P'\) & Condition at site \(i\)\\
\midrule
\(J_{3,L}\) &
\(\sigma\in\{X_i,Y_i\},\ P'\in\mathcal S_3\) & \(P'_i=Z\)\\
\(J_{3,Z}\) &
\(\sigma=Z_i,\ P'\in\mathcal S_3\) & \(P'_i\in\{X,Y\}\)\\
\(J_{3,\perp}\) &
\(\sigma\in\{X_i,Y_i\},\ P'\in\mathcal S_3\) &
\(P'_i\in\{X,Y\},\ P'_i\ne\sigma_i\)\\
\(J_{1,\perp}\) &
\(\sigma\in\{X_i,Y_i\},\ P'\in\mathcal S_1\) &
\(P'_i\in\{X,Y\},\ P'_i\ne\sigma_i\)\\
\(J_{1,Z}\) &
\(\sigma=Z_i,\ P'\in\mathcal S_1\) & \(P'_i\in\{X,Y\}\)\\
\bottomrule
\end{tabular}
\end{center}
All five sums vanish if the common-antiunitary defects vanish.

The certificate polynomial is
\begin{align}
 \mathcal Q(T)={}&-\sum_{E\in\mathcal V_3}c_E^2
 +\frac5{12}J_{3,L}+\frac56J_{3,Z}
 +\frac7{12}J_{3,\perp}+\frac12g\nonumber\\
 &+\frac32J_{1,\perp}-\frac32J_{1,Z}
 +\frac{17}{6}f_{\mathcal F_4}^{13}
 +\frac{19}{6}f_{\mathcal F_4}^{33}
 +\frac{25}{6}f_{\mathcal D_3}^{13}
 +\frac13f_{\mathcal D_2}^{33}.
 \label{o2gc:Q}
\end{align}

\begin{lemma}[Exact coefficient matrix and positivity]
\label{o2gc:lem-matrix}
Order the 32 Pauli labels in \(\mathcal S_1\cup\mathcal S_3\)
lexicographically using \(I<X<Y<Z\). Let \(z\) contain all
\(528\) monomials \(t_{P'}t_{Q'}\), \(P'\le Q'\), with no
\(\sqrt2\) factors on unequal pairs. There is an integer real
symmetric \(528\times528\) matrix \(M\) such that, identically for
every complex \(T\) in this coordinate space,
\begin{equation}
 \mathcal Q(T)=\frac1{12}z^\dagger Mz.
 \label{o2gc:matrix-identity}
\end{equation}
The \(36\) rows and columns indexed by two \(\mathcal S_1\)
coordinates are zero. On the remaining \(492\) coordinates,
\(M\) has \(111\) blocks and satisfies
\begin{equation}
 M_{\mathrm{remaining}}-4I_{492}>0.
 \label{o2gc:matrix-positive}
\end{equation}
In particular,
\begin{equation}
 \mathcal Q(T)\ge\frac13
 \sum_{\substack{P'\le Q'\\
 P'\in\mathcal S_3\ \text{or}\ Q'\in\mathcal S_3}}
 |t_{P'}t_{Q'}|^2.
 \label{o2gc:coercive}
\end{equation}
\end{lemma}

\begin{proof}[Exact finite verification]
Treat the \(t_{P'}\) and their formal conjugates as independent
variables and expand \(12\mathcal Q\) in bidegree \((2,2)\).
Each monomial has the unique form
\(\overline{t_{P'}t_{Q'}}t_{R'}t_{S'}\), with \(P'\le Q'\) and
\(R'\le S'\). Accumulating all ordered trace contributions into
these unordered pair indices uniquely determines the coefficient
matrix \(M\); it is not a fitted Gram representation.

The exact verifier multiplies Paulis with phases \(1,-1,i,-i\),
expands the normalized traces and real parts in
\eqref{o2gc:scalar-contractions}--\eqref{o2gc:Q}, and checks integral
entries and \(M=M^T\). In particular it proves
\eqref{o2gc:matrix-identity} before imposing any encoder equations.
No floating-point eigenvalues or numerical matrix entries enter
this verification.

Conjugation by a four-qubit Pauli multiplies each \(t_{P'}\) by its
Pauli-conjugation sign and leaves \(\mathcal Q\) invariant. Hence a
matrix entry between the monomials \(t_{P'}t_{Q'}\) and
\(t_{R'}t_{S'}\) can be nonzero only when the Pauli labels of
\(P'Q'\) and \(R'S'\) agree modulo phase. This gives the character
blocks. The verifier checks every off-block zero entry, all \(36\)
zero rows and columns, and every remaining block.

The exact characteristic polynomials of all blocks are summarized
in Table~\ref{o2gc:tab-blocks}. Every linear root is greater than
four. Each quadratic factor \(x^2-bx+c\) has real roots because it
is a factor of the characteristic polynomial of a real symmetric
matrix. Its shifted roots \(y=x-4\) have positive sum \(b-8\) and
positive product \(c-4b+16\), as is checked from every listed factor.
Thus both roots exceed four. This proves
\eqref{o2gc:matrix-positive} once the exact block polynomials have
been verified.

The standard-library verifier also establishes positivity without
factoring characteristic polynomials. For each of all \(111\)
integer blocks, it subtracts \(4I\) and performs exact rational
\(LDL^T\) elimination. Every pivot is strictly positive. Explicitly,
at pivot \(d>0\), the remaining Schur complement is updated by
\[
 A_{ij}\longmapsto A_{ij}-\frac{A_{ik}A_{kj}}{d}.
\]
All entries and pivots are rational, and positive pivots give positive
definiteness by congruence. Every block is checked individually, so
this argument does not assume that one representative block covers
its entire symmetry class.

Finally, the zero coordinates are precisely the
\(\binom{8+1}{2}=36\) products of two weight-one coordinates.
On the remaining coordinates \eqref{o2gc:matrix-positive} implies
\eqref{o2gc:coercive}. The matrix is defined completely by
\eqref{o2gc:Q} and the specified monomial ordering; the accompanying
JSON additionally records every block and every positive pivot.
\end{proof}

\begin{table}[htbp]
\centering
\small
\renewcommand{\arraystretch}{1.2}
\begin{tabular}{@{}lrrp{0.64\linewidth}@{}}
\toprule
Character pattern & Size & Count & Characteristic polynomial of a block of \(M\)\\
\midrule
\(I\)&24&1&
\((x-424)(x-192)^4(x-24)^{11}(x-16)^2
(x^2-48x+192)^3\)\\
\(Z\)&3&4&\((x-384)(x-48)^2\)\\
\(pq\)&3&24&\((x-52)(x^2-460x+47872)\)\\
\(pqZ\)&4&48&\((x-50)(x-46)(x^2-248x+10716)\)\\
\(ZZ\)&8&6&
\((x-32)(x-24)(x^2-408x+9856)(x^2-352x+9216)^2\)\\
\(ZZZ\)&12&4&
\((x-136)(x-48)(x^2-352x+9216)^3(x^2-64x+720)^2\)\\
\(pqZZ\)&4&24&\((x-56)(x-44)(x^2-380x+23936)\)\\
\bottomrule
\end{tabular}
\caption{All exact character blocks after removal of the \(36\) zero
\(\mathcal S_1\mathcal S_1\) coordinates. Here \(p,q\in\{X,Y\}\);
all distinct-site and transverse-axis choices are included.
The counts sum to \(111\), and the sizes with multiplicity sum to
\(492\). Each displayed polynomial is verified for every block in
the corresponding class.}
\label{o2gc:tab-blocks}
\end{table}

\begin{proof}[Proof of Proposition~\ref{o2gc:prop-rigidity}]
Under \eqref{o2gc:rigidity-hypotheses}, \(g\) and all five \(J\)-sums
vanish. Each forbidden-set contraction also vanishes because its
factor \(c_E\) is zero. Hence \eqref{o2gc:Q} and
\eqref{o2gc:coercive} give
\[
 -\sum_{E\in\mathcal V_3}c_E^2
 =\mathcal Q(T)
 \ge\frac13
 \sum_{\substack{P'\le Q'\\
 P'\in\mathcal S_3\ \text{or}\ Q'\in\mathcal S_3}}
 |t_{P'}t_{Q'}|^2\ge0.
\]
The left side is nonpositive, so every \(c_E\),
\(E\in\mathcal V_3\), vanishes. The right side includes
\(|t_{P'}|^4\) for every \(P'\in\mathcal S_3\), so all those
coordinates vanish as well. Thus \(T_3=0\).
\end{proof}

\subsection{Completion of the converse and exact attainment}
\label{o2gc:subsec-finish}

\begin{proof}[Proof of Theorem~\ref{o2gc:thm-image}]
For every exact orbit-2 encoder, the analytic construction gives
\eqref{o2gc:T-support}, the partial-isometry identities, and all twelve
physical equations \eqref{o2gc:physical-J}. As noted after
Proposition~\ref{o2gc:prop-rigidity}, these imply every hypothesis
\eqref{o2gc:rigidity-hypotheses}. The proposition therefore gives
\(c_E=0\) for \(E\in\mathcal V_3\). By
\eqref{o2gc:normalization}, the coefficient of \(E\) in
\(K=2TT^\dagger-I\) is \(2c_E\), so
\begin{equation}
 A_3=\frac1{16}\operatorname{Tr}(K_3^2)
     =4\sum_{E\in\mathcal V_3}c_E^2=0.
 \label{o2gc:A3-zero}
\end{equation}
Equation~\eqref{o2gc:plane} gives \(\mathcal L(\alpha)\), with
\(0\le\alpha\le1\) by \(A_1,A_2\ge0\).

For completeness, the reverse inclusion has the following explicit
split-qubit construction. Start with the standard perfect five-qubit
code with stabilizer generators
\[
 XZZXI,\quad IXZZX,\quad XIXZZ,\quad ZXIXZ.
\]
It is pure: every nonidentity Pauli of weight at most two has zero
compression. Its logical \(Z^{\otimes5}\), multiplied by the first
generator, gives the weight-three logical representative \(-YIIYZ\).
Local Clifford rotations and a permutation therefore give a pure
five-qubit code with a chosen logical \(Z_1Z_2Z_3\).

For \(0\le\theta\le\pi/2\), put \(a=\cos(\theta/2)\),
\(b=\sin(\theta/2)\), and split original site \(1\) by the isometry
\[
 W_\theta\ket0=a\ket{00}+b\ket{11},\qquad
 W_\theta\ket1=a\ket{01}+b\ket{10}.
\]
Its exact identities are
\begin{equation}
 W_\theta W_\theta^\dagger
 =\frac12(I+\cos\theta\,Z\otimes I+\sin\theta\,X\otimes X),
 \qquad (Z\otimes Z)W_\theta=W_\theta Z.
 \label{o2gc:split-construction}
\end{equation}
Take the two split outputs as \(A_1,A_2\), original sites \(2,3\)
as \(A_3,A_4\), and original sites \(4,5\) as \(B_1,B_2\).
The original weight-three logical lifts to \(I_{\mathsf B}D\),
which is precisely \eqref{o2gc:record}.

Compression of any six-qubit error of weight at most two through
\(W_\theta\otimes I^{\otimes4}\) is supported on at most two original
sites. Purity of the five-qubit code therefore proves all KL equations.
If the error has a nonidentity factor on an unsplit site, its scalar
compression is zero. On the split pair its scalar is the identity
coefficient \(\operatorname{Tr}(W_\theta^\dagger E W_\theta)/2\).
Consequently the only nonzero weight-one and weight-two coefficients
are \(\cos\theta\) on \(Z_{A_1}\) and \(\sin\theta\) on
\(X_{A_1}X_{A_2}\). Thus
\[
 A_1=\cos^2\theta=\alpha,\qquad A_2=\sin^2\theta=1-\alpha.
\]
The universal inequality \eqref{o2gc:shadow} gives \(A_3=0\), and
\eqref{o2gc:MacWilliams} gives \(\mathcal L(\alpha)\). Every
\(\alpha\in[0,1]\), including both endpoints, occurs. The quantum
Singleton bound excludes distance at least four at length six and
code dimension two, so these are exact \(((6,2,3))\) realizations.
\end{proof}

\begin{remark}[Scope of the converse]
The proof uses no division by a computational coefficient, row norm,
polarization, Schmidt eigenvalue, or Gram minor. It covers arbitrary
interactions among different \(Z\)-pairs in \(T_3\) and both regular
and singular frame configurations. The necessary common-antiunitary
equations already suffice for this enumerator conclusion; no
sufficiency claim for the matrix completion problem is required.
The theorem determines \(\mathcal W_2\), not all code projectors up to
local unitaries, and does not assert global connectedness or uniqueness
of the split-code family. It also does not prove a stronger
partial-isometry inequality with the common-antiunitary equations
omitted.
\end{remark}

\subsection{Exact certificate files and reproducibility}
\label{o2gc:subsec-reproducibility}

The proof and its exact certificate are distributed together in
\path{orbit2_enumerator_package.zip}. All filenames below are relative
to the root of the unpacked package:
\begin{itemize}
\item \path{orbit2_enumerator_section.tex}:
this include-ready section, using the supplement's existing preamble;
\item \path{orbit2_enumerator_section_preview.tex}:
a standalone document containing the same section and its full preamble;
\item \path{orbit2_exact_certificate.py}:
the standard-library verifier that regenerates the integer matrix
from the quartic polynomial and checks all rational pivots;
\item \path{orbit2_certificate_blocks.json}:
all integer blocks, monomial indices, and positive rational pivots;
\item \path{orbit2_verification_log.txt}:
the successful exact-verification record regenerated for this package;
\item \path{README.txt}:
instructions for compiling, including, and verifying the proof;
\item \path{SHA256SUMS}:
SHA-256 file hashes for both LaTeX sources, the verifier, the certificate
JSON, the verification log, and the README.
\end{itemize}
Run from the unpacked package directory with Python~3, without
third-party packages:
\begin{verbatim}
python3 orbit2_exact_certificate.py
python3 orbit2_exact_certificate.py --json regenerated_blocks.json
\end{verbatim}
The regenerated JSON agrees with the supplied
\path{orbit2_certificate_blocks.json}. The verifier reconstructs the
matrix directly from the explicitly coded polynomial identity; it does
not read the LaTeX files or trust a precomputed positive-semidefinite
matrix. The JSON and log are supporting records rather than runtime
dependencies.
The canonical sparse integer matrix has \(2{,}844\) nonzero entries.
Its SHA-256 under the verifier's specified sparse-entry serialization is
\begin{center}
\small\ttfamily
64b192ccce90edf4b583285307f46265435a51b032f72e6493057138a1820b1d
\end{center}
Discovery of the multipliers can use numerical optimization; the proof
does not. Its finite verification uses exact integer coefficient
expansion and rational \texttt{Fraction} arithmetic. The analytic
encoder reduction is proved above, so this section does not depend
on an unproved converse or on a regular-support assumption from an
earlier problem note.

% This file has no document class, preamble, or document environment.
% It uses only the supplement's existing packages, macros, and theorem styles.
% Every new label has the prefix o2gc:.
% Existing abstract/status/count statements elsewhere need separate updating;
% this fragment does not edit or silently replace those statements.

\section{Exact weight-enumerator image of orbit 2}
\label{o2gc:sec}

Orbit~2 is the primitive order-two record
\[
 (q,m,a)=\bigl(2,2,(0,0,1,1,1,1)\bigr).
\]
For the proof, relabel the physical sites in the order
\(B_1,B_2,A_1,A_2,A_3,A_4\). Write
\(\mathsf B=\{B_1,B_2\}\) for the inactive pair and
\(\mathsf A=\{A_1,A_2,A_3,A_4\}\) for the active sites. Its chosen physical
logical involution is
\begin{equation}
 D_{\mathrm{phys}}=D(a;2)=I_{\mathsf B}\otimes D,
 \qquad D=Z_{A_1}Z_{A_2}Z_{A_3}Z_{A_4}.
 \label{o2gc:record}
\end{equation}
Let \(V:\C^2\to(\C^2)^{\otimes6}\) be an isometric encoder, with
code projector \(P=VV^\dagger\), satisfying
\begin{equation}
 D_{\mathrm{phys}}V=VZ_L,\qquad
 V^\dagger EV=\lambda_EI_2\quad(\wt(E)\le2),
 \label{o2gc:KL}
\end{equation}
where \(E\) ranges over Hermitian Pauli strings. For all Pauli weights
put \(\lambda_E=\operatorname{Tr}(PE)/2\); this agrees with the KL scalar
at weights at most two but does not assert scalar compression at higher
weights. We use the supplement's normalization
\[
 A_w=\frac14\sum_{\wt(E)=w}|\operatorname{Tr}(PE)|^2
     =\sum_{\wt(E)=w}\lambda_E^2.
\]

\begin{theorem}[Complete orbit-2 enumerator image]
\label{o2gc:thm-image}
Every exact orbit-2 realization has \(A_3=0\). Its complete
Shor--Laflamme \(A\)-enumerator image is
\begin{equation}
 \boxed{
 \mathcal W_2=
 \left\{\mathcal L(\alpha):0\le\alpha\le1\right\},
 \qquad
 \mathcal L(\alpha)=
 (1,\alpha,1-\alpha,0,11+4\alpha,16-\alpha,3-3\alpha).
 }
 \label{o2gc:image}
\end{equation}
The converse includes every zero-support and singular stratum. No
stabilizer assumption is made.
\end{theorem}

The proof has an analytic reduction and an exact finite certificate.
First, KL and complementary purities force a flat inactive pair and a
specific rank-eight active projector. Mixed KL supplies one common
antiunitary, whose spectral saturation restricts an associated partial
isometry to 32 complex Pauli coordinates. A homogeneous quartic
certificate then forces its weight-three coordinates, and the code's
entire \(A_3\), to vanish. The certificate identity and positivity are
verified with integer and rational arithmetic.

\subsection{Flatness of the inactive pair and the active reflection}
\label{o2gc:subsec-flatness}

Put \(\rho=P/2\). The universal six-qubit rank-two MacWilliams normal
form, with \(x=A_1,y=A_2,z=A_3\), is
\begin{equation}
 (A_0,\ldots,A_6)=
 (1,x,y,z,15-4y-3z,12+3x+4y+3z,4-4x-y-z).
 \label{o2gc:MacWilliams}
\end{equation}
It follows from the quantum MacWilliams transform with factor \(1/32\)
and \(A_j=B_j\) for \(j=0,1,2\). With full six-qubit spin flip
\(\widetilde P=Y^{\otimes6}P^TY^{\otimes6}\), positivity gives
\begin{equation}
 0\le\operatorname{Tr}(P\widetilde P)
 =\frac1{16}\sum_{w=0}^6(-1)^wA_w
 =\frac12(1-A_1-A_2-A_3).
 \label{o2gc:shadow}
\end{equation}

\begin{lemma}[Universal flatness and weight-three support]
\label{o2gc:lem-flatness}
Every encoder in \eqref{o2gc:KL} satisfies
\begin{equation}
 \rho_{\mathsf B}=\frac{I_4}{4},\qquad
 \rho_{\mathsf B A_i}=\frac{I_4}{4}\otimes\rho_{A_i}
 \quad(1\le i\le4),\qquad A_1+A_2+A_3=1.
 \label{o2gc:flatness}
\end{equation}
The only possibly nonzero weight-three coefficients are
\(\lambda_{p_iq_jZ_k}\) on distinct active sites, with
\(p,q\in\{X,Y\}\). Moreover,
\begin{equation}
 \rho_{\mathsf A}=\frac{\Pi}{8},\qquad
 \Pi^\dagger=\Pi=\Pi^2,\qquad \rank\Pi=8,\qquad [\Pi,D]=0,
 \label{o2gc:active-projector}
\end{equation}
and each active parity block of \(\Pi\) has rank four.
\end{lemma}

\begin{proof}
Purify the maximally mixed code with a reference qubit \(R\):
\[
 \ket{\Psi}=
 \frac{\ket0_R\ket{0_L}+\ket1_R\ket{1_L}}{\sqrt2}.
\]
The operator \(G=Z_RI_{\mathsf B}D\) fixes this state. KL says that every
correlation involving a nonidentity reference Pauli and a physical
Pauli of weight at most two vanishes. Multiplication by \(G\) preserves
squared correlations.

Physical coefficients with an odd number of active \(X/Y\) factors
vanish by \(D\)-parity. An active-only coefficient with at least two
\(Z\) factors also vanishes: multiplying by \(G\) gives a
reference-\(Z\) correlation with physical weight \(4-z\le2\), where
\(z\) is the original active \(Z\)-count.
For this purity computation, denote the remaining nonidentity
active-only squared masses by
\[
 \begin{array}{c|cccc}
 \text{mass}&a&u&v&w\\ \hline
 \text{pattern}&Z_i&p_iq_j&p_iq_jZ_k&p_1p_2p_3p_4
 \end{array}
 \qquad(p,q,p_i\in\{X,Y\}).
\]
All displayed sites in a pattern are distinct. Let \(s=a+u+v+w\), and
define the nonnegative masses
\[
 b=\sum_{\substack{F\in\{I,X,Y,Z\}^{\otimes\mathsf B}\\F\ne I}}
       \lambda_F^2,\qquad
 m=\sum_{i=1}^4\sum_{F\ne I}\lambda_{FZ_i}^2.
\]
Here \(F\) in the second sum ranges over all nonidentity spectator
Paulis, including those of spectator weight two.

Purity of a marginal on \(r\) qubits is \(2^{-r}\) times the sum of its
squared Pauli correlations. KL gives
\(\rho_{R\mathsf B}=I_R/2\otimes\rho_{\mathsf B}\).
Equality of complementary purities for \(\mathsf A\) and \(R\mathsf B\)
therefore gives
\begin{equation}
 1+s=2(1+b).
 \label{o2gc:purity-1}
\end{equation}
Let \(h\) be the squared correlation mass on \(R\mathsf A\) with
reference factor \(X\) or \(Y\). Multiplication by \(G\) identifies the
reference-\(Z\) mass with the physical active mass \(1+s\). Comparing
the complementary purities of \(\mathsf B\) and \(R\mathsf A\) gives
\[
 2(1+s)+h=8(1+b),\qquad h=4(1+b).
\]
A surviving reference-\(X/Y\) correlation has odd active transverse
count. If that count is one and its active \(Z\)-count is \(z\), the
term and its \(G\)-multiple have physical weights \(1+z\) and \(4-z\).
At least one weight is at most two, so both correlations vanish.
Thus the surviving transverse count is three. Multiplication by \(G\)
interchanges its physical-weight-three and physical-weight-four
terms, giving
\begin{equation}
 h_3=h_4=\frac h2=2(1+b).
 \label{o2gc:purity-h}
\end{equation}

For each \(i\), compare the complementary marginals
\(\mathsf B A_i\) and \(R(\mathsf A\setminus\{A_i\})\), then sum over
all four active sites. The first marginal's mass is
\[
 1+b+\lambda_{Z_i}^2+\sum_{F\ne I}\lambda_{FZ_i}^2.
\]
On the complementary side, reference identity and reference \(Z\)
together count active strings whose \(i\)th factor is \(I\) or \(Z\).
A string with transverse count \(t\) is therefore counted \(4-t\)
times. Reference-\(X/Y\) physical-weight-three terms appear once and
weight-four terms never appear. Accounting also for the dimensions of
the two marginals gives
\begin{equation}
 2\{4(1+b)+a+m\}
 =4(1+s)-2u-2v-4w+h_3.
 \label{o2gc:purity-2}
\end{equation}
Substitution of \eqref{o2gc:purity-1} and \eqref{o2gc:purity-h} yields
\begin{equation}
 b+m+w=0.
 \label{o2gc:purity-sos}
\end{equation}
Every summand is a sum of real squares. Hence \(b=m=w=0\), every
coefficient they contain vanishes, and \(a+u+v=1\). In particular
\(\rho_{\mathsf B}=I_4/4\). Parity kills spectator-dressed active
\(X/Y\) coefficients, and \(m=0\) kills all spectator-dressed \(Z_i\)
coefficients. This proves the product marginal in
\eqref{o2gc:flatness}.

The full weight-one and weight-two masses are now \(A_1=a\) and
\(A_2=u\). Write \(A_3=v+n\), where \(n\ge0\) collects the remaining
terms supported on two active sites and one spectator site. All other
additional low-weight patterns have vanished. Equation
\eqref{o2gc:shadow} gives \(1+n\le1\), so \(n=0\), \(A_3=v\), and
\(A_1+A_2+A_3=1\).

Finally, KL for all errors on \(\mathsf B\) gives the same marginal
\(I_4/4\) for each logical word. Each word therefore has four equal
Schmidt coefficients \(1/2\). Its active Schmidt space lies in the
corresponding \(D\)-parity sector. Each logical word's active marginal
is its rank-four projector divided by four. Averaging these two
marginals gives \(\Pi/8\), proving
\eqref{o2gc:active-projector}.
\end{proof}

Write the resulting flat parity frames as \(16\times4\) matrices
\(W_0,W_1\):
\begin{equation}
 \ket{z_L}=\frac12\sum_{b=0}^3\ket b_{\mathsf B}\otimes W_z\ket b,
 \qquad W_z^\dagger W_z=I_4,\qquad DW_z=(-1)^zW_z.
 \label{o2gc:frames}
\end{equation}
Their ranges are orthogonal, and
\[
 \Pi=W_0W_0^\dagger+W_1W_1^\dagger
     =4\operatorname{Tr}_{\mathsf B}P.
\]
Set \(K=2\Pi-I_{16}\). Then \(K^\dagger=K\), \(K^2=I_{16}\), and
\([K,D]=0\). For every nonidentity active Pauli \(E\),
\begin{equation}
 \lambda_{I_{\mathsf B}E}
 =\frac18\operatorname{Tr}(\Pi E)
 =\frac1{16}\operatorname{Tr}(KE).
 \label{o2gc:normalization}
\end{equation}
The active support proved above thus gives
\begin{align}
 K&=K_1+K_2+K_3,\label{o2gc:K-support}\\
 K_1&=\sum_{i=1}^4a_iZ_i,\nonumber\\
 K_2&=\sum_{i<j}\sum_{p,q\in\{X,Y\}}u_{ij}^{pq}p_iq_j,\nonumber\\
 K_3&=\sum_{i<j}\sum_{k\notin\{i,j\}}
              \sum_{p,q\in\{X,Y\}}v_{ij;k}^{pq}p_iq_jZ_k.\nonumber
\end{align}
The three sets have \(4,24,48\) real coefficients, respectively. Their
squared masses are
\begin{equation}
 \alpha=\frac1{16}\operatorname{Tr}(K_1^2)=A_1,\qquad
 \beta=\frac1{16}\operatorname{Tr}(K_2^2)=A_2,\qquad
 \nu=\frac1{16}\operatorname{Tr}(K_3^2)=A_3,
 \quad \alpha+\beta+\nu=1.
 \label{o2gc:masses}
\end{equation}
Consequently \eqref{o2gc:MacWilliams} already reduces the enumerator to
\begin{equation}
 (A_0,\ldots,A_6)=
 (1,\alpha,1-\alpha-\nu,\nu,
   11+4\alpha+\nu,16-\alpha-\nu,3-3\alpha).
 \label{o2gc:plane}
\end{equation}
The remaining converse is \(\nu=0\). An arbitrary reflection having
\eqref{o2gc:K-support} need not be completable to a six-qubit code; the
mixed KL constraints are essential in the next step.

\subsection{One common antiunitary from mixed KL}
\label{o2gc:subsec-commonJ}

On the spectator pair choose the fixed basis of
\(\Theta_{\mathsf B}=(Y\otimes Y)\mathcal K\), where \(\mathcal K\)
denotes computational complex conjugation. Explicitly, let \(M_{\mathsf B}\)
have columns
\[
 i\Phi_+,\quad\Phi_-,\quad\Psi_+,\quad i\Psi_-,
 \qquad
 \Phi_\pm=\frac{\ket{00}\pm\ket{11}}{\sqrt2},\quad
 \Psi_\pm=\frac{\ket{01}\pm\ket{10}}{\sqrt2}.
\]
Each column is fixed by \(\Theta_{\mathsf B}\). This is a common
spectator magic basis, not an independently chosen basis for each
observable. In that basis put
\[
 Q_z=W_z\overline{M_{\mathsf B}},\qquad Q=(Q_0\ Q_1).
\]
The conjugate follows from the bipartite expansion
\eqref{o2gc:frames}. We have \(Q^\dagger Q=I_8\) and \(QQ^\dagger=\Pi\).

\begin{lemma}[Common antiunitary equations]
\label{o2gc:lem-commonJ}
For every real single-body active Hamiltonian \(H\), put
\[
 h=Q^\dagger HQ,\qquad
 \mu_H=\frac18\operatorname{Tr}(\Pi H).
\]
Then the same skew-unitary matrix
\begin{equation}
 J_0=\begin{pmatrix}0&I_4\\-I_4&0\end{pmatrix}
 \label{o2gc:J0}
\end{equation}
satisfies
\begin{equation}
 hJ_0+J_0h^T=2\mu_HJ_0.
 \label{o2gc:commonJ}
\end{equation}
\end{lemma}

\begin{proof}
Each single-spectator Pauli \(F\) changes sign under
\(\Theta_{\mathsf B}\). Hence \(M_{\mathsf B}^\dagger F M_{\mathsf B}\)
is imaginary, antisymmetric, and Hermitian. The six such Paulis span
the six-dimensional space of imaginary antisymmetric Hermitian
\(4\times4\) matrices.

For every single-active Pauli \(E\), mixed KL gives, for all six \(F\)
and all \(z,z'\),
\[
 \frac14\operatorname{Tr}\!\left[
 (M_{\mathsf B}^\dagger F M_{\mathsf B})^T
 Q_z^\dagger E Q_{z'}\right]=0.
\]
Orthogonality to this antisymmetric space makes every block
\(Q_z^\dagger EQ_{z'}\) complex symmetric. A diagonal block is also
Hermitian and is therefore real symmetric.

The stronger product-marginal identity in \eqref{o2gc:flatness}, not
merely spectator flatness, gives
\[
 h_{00}+h_{11}=2\mu_HI_4.
\]
Indeed, tracing \((I_{\mathsf B}\otimes H)\rho\) over the active sites
gives the spectator matrix \((h_{00}+h_{11})^T/8\) in this basis;
the product marginal instead gives \(\mu_HI_4/4\).
Thus
\[
 h=\begin{pmatrix}
 R&C\\ \overline C&2\mu_HI_4-R
 \end{pmatrix},
 \qquad R=R^T=\overline R,\qquad C=C^T.
\]
Direct block multiplication proves \eqref{o2gc:commonJ}.
\end{proof}

The antiunitary \(J_0\mathcal K\) squares to \(-I_8\) and reverses
every centered compression \(h-\mu_HI_8\). The same \(J_0\) works for
all twelve single-active Paulis. Separate spectral pairings, or a
condition using only the four \(Z_i\), would not provide the full
hypotheses used in the final certificate.

\subsection{A sparse partial isometry on four active qubits}
\label{o2gc:subsec-T}

Define
\begin{equation}
 \Omega=QJ_0Q^T=Q_0Q_1^T-Q_1Q_0^T,\qquad
 S=Y^{\otimes4},\qquad T=\Omega S.
 \label{o2gc:T-definition}
\end{equation}
Here \(S\) is real, symmetric, and unitary. Using
\(Q^T\overline Q=I_8\) gives
\[
 \Omega^T=-\Omega,\qquad
 \Omega\Omega^\dagger=\Pi,\qquad
 \Omega^\dagger\Omega=\Pi^T.
\]
Therefore
\begin{equation}
 TT^\dagger=\Pi,\qquad
 T^\dagger T=\widetilde\Pi:=S\Pi^TS,\qquad
 ST^TS=-T,\qquad \operatorname{Tr}(TT^\dagger)=8.
 \label{o2gc:partial-isometry}
\end{equation}
The Pauli identity \(SE^TS=(-1)^{\wt(E)}E\) means that only odd
Pauli weights can occur in \(T\).

\begin{lemma}[Exact edge sparsity]
\label{o2gc:lem-edges}
For computational active strings \(x,y\in\{0,1\}^4\),
\[
 d_H(x,y)=1\quad\Longrightarrow\quad\Omega_{xy}=0.
\]
In fact \(\Omega\) is supported only at Hamming distance three, and
\(T\) only at Hamming distance one.
\end{lemma}

\begin{proof}
Parity already makes \(\Omega\) zero between equal-parity strings,
since it pairs \(Q_0\) and \(Q_1\). Fix an even string \(x\), set
\(s_i=(-1)^{x_i}\), \(t=\sum_i a_is_i\), and write \(\bar x\) for
its complement. Consider
\[
 H_x=\sum_i s_iZ_i,\qquad
 r_x=\bra xQ_0,\qquad r_{\bar x}=\bra{\bar x}Q_0.
\]
The allowed reflection support \eqref{o2gc:K-support} gives
\[
 \|r_x\|^2=\Pi_{xx}=\frac{1+t}{2},\qquad
 \|r_{\bar x}\|^2=\frac{1-t}{2},\qquad
 r_xr_{\bar x}^\dagger=\Pi_{x\bar x}=0.
\]
The last equality uses the absence of any term flipping all four
active bits.

The eigenvalue of \(H_x\) on \(y\) is \(4-2d_H(x,y)\). Thus its
even-parity restriction is
\(4(\ket x\bra x-\ket{\bar x}\bra{\bar x})\), and its odd-parity
spectrum is \(\{-2,+2\}\). With \(h_z=Q_z^\dagger H_xQ_z\), this gives
\[
 h_0=4r_x^\dagger r_x-4r_{\bar x}^\dagger r_{\bar x},\qquad
 h_0^Tr_x^T=2(1+t)r_x^T.
\]
Moreover \(\mu_{H_x}=t\). Equation~\eqref{o2gc:commonJ} implies
\[
 h_1=2tI_4-h_0^T,\qquad h_1r_x^T=-2r_x^T.
\]
Set \(w=Q_1r_x^T\). Then
\[
 \bra w(H_x+2I)\ket w=0.
\]
On odd parity \(H_x+2I\) is positive semidefinite, with kernel
consisting exactly of the strings at distance three from \(x\).
Consequently \(w\) has no distance-one components. For even \(x\)
and odd \(y\),
\[
 \Omega_{xy}=r_x(\bra yQ_1)^T=\bra y w,
\]
so the desired entries vanish. If \(r_x=0\), that row already vanishes;
no division is needed. Skew symmetry supplies the reverse ordering.
Only odd distances are possible, so the remaining distance is three.
Right multiplication by \(S\) complements the column string up to a
phase, taking distance three to distance one.
\end{proof}

An operator supported only on distance-one matrix entries has exactly
one \(X/Y\) factor in its Pauli expansion. Combining this with the odd
Pauli weights in \eqref{o2gc:partial-isometry} gives
\begin{align}
 T&=T_1+T_3,\label{o2gc:T-support}\\
 T_1&=\sum_i(t_i^XX_i+t_i^YY_i),\nonumber\\
 T_3&=\sum_i\sum_{\substack{j<k\\j,k\ne i}}
          (t_{i;jk}^XX_i+t_{i;jk}^YY_i)Z_jZ_k.\nonumber
\end{align}
There are \(8+24=32\) complex coefficients.

For each single-active Pauli \(\sigma\), multiply
\eqref{o2gc:commonJ} by \(Q,Q^T\) to obtain
\[
 \Pi\sigma\Omega+\Omega\sigma^T\Pi^T=2\mu_\sigma\Omega.
\]
Since \(S\sigma^TS=-\sigma\), multiplication on the right by \(S\)
yields the physical common-antiunitary equation
\begin{equation}
 \Pi\sigma T-T\sigma\widetilde\Pi=2\mu_\sigma T,\qquad
 \mu_\sigma=\frac18\operatorname{Tr}(\Pi\sigma).
 \label{o2gc:physical-J}
\end{equation}
All steps so far are valid for arbitrary computational zeros and
singular overlap patterns.

\subsection{A homogeneous rigidity theorem}
\label{o2gc:subsec-rigidity}

The remaining finite statement can be formulated independently of
encoders. All operators in this subsection act on the four active
qubits; active identity factors are suppressed. Write
\[
 \tau(M)=\frac1{16}\operatorname{Tr}(M),\qquad
 \mathcal S_1=\{X_i,Y_i:1\le i\le4\},
\]
\[
 \mathcal S_3=\{p_iZ_jZ_k:p\in\{X,Y\},\ j<k,\ i\notin\{j,k\}\}.
\]
For an arbitrary \(T\) in the complex span of
\(\mathcal S_1\cup\mathcal S_3\), let \(T_1,T_3\) denote these two
coordinate parts and put
\begin{equation}
 N=\tau(TT^\dagger),\qquad c_E=\tau(ETT^\dagger).
 \label{o2gc:homogeneous-scalars}
\end{equation}
Each \(c_E\) is real. Define the homogeneous cubic defects
\begin{align}
 F(T)&=TT^\dagger T-2NT,\label{o2gc:defects}\\
 R_\sigma(T)&=TT^\dagger\sigma T-T\sigma T^\dagger T-4c_\sigma T.
 \nonumber
\end{align}
The target and forbidden Pauli sets are
\begin{align}
 \mathcal V_3&=\{p_iq_jZ_k:i<j,\ k\notin\{i,j\},\
                                  p,q\in\{X,Y\}\},\nonumber\\
 \mathcal D_2&=\{Z_iZ_j:i<j\},\nonumber\\
 \mathcal D_3&=\{Z_iZ_jZ_k:i<j<k\},\nonumber\\
 \mathcal F_4&=\{p_iq_jZ_kZ_\ell:i<j,\
                \{k,\ell\}=\{1,2,3,4\}\setminus\{i,j\},\
                p,q\in\{X,Y\}\}.
 \label{o2gc:Pauli-sets}
\end{align}
Their sizes are \(48,6,4,24\), respectively.

\begin{proposition}[Global common-antiunitary rigidity]
\label{o2gc:prop-rigidity}
Suppose \(T\) has the support \eqref{o2gc:T-support} and
\begin{equation}
 F(T)=0,\qquad R_\sigma(T)=0
 \quad\text{for all twelve single-active Paulis},\qquad
 c_E=0\quad(E\in\mathcal D_2\cup\mathcal D_3\cup\mathcal F_4).
 \label{o2gc:rigidity-hypotheses}
\end{equation}
Then \(T_3=0\) and \(c_E=0\) for every \(E\in\mathcal V_3\).
\end{proposition}

For the matrix constructed from an orbit-2 code,
\(N=\tau(\Pi)=1/2\), so \(F(T)=0\) follows from its partial-isometry
identities. Also \(c_\sigma=\mu_\sigma/2\), so
\eqref{o2gc:physical-J} gives \(R_\sigma(T)=0\) with exactly the
factor \(4\) in \eqref{o2gc:defects}. Finally, the support of
\(K=2TT^\dagger-I\) gives all the forbidden zeros in
\eqref{o2gc:rigidity-hypotheses}. Thus the proposition applies once
its certificate is proved.

\subsection{The exact positive quartic certificate}
\label{o2gc:subsec-certificate}

Define the scalar defect contractions
\begin{align}
 g&=\Re\,\tau(T_3^\dagger F(T)),\nonumber\\
 c_E^{13}&=\tau\!\left(E(T_1T_3^\dagger+T_3T_1^\dagger)\right),
 \qquad c_E^{33}=\tau(ET_3T_3^\dagger),\nonumber\\
 f_{\mathcal C}^{13}&=\sum_{E\in\mathcal C}c_Ec_E^{13},\qquad
 f_{\mathcal C}^{33}=\sum_{E\in\mathcal C}c_Ec_E^{33}.
 \label{o2gc:scalar-contractions}
\end{align}
The last two definitions apply to any Pauli set \(\mathcal C\).
All displayed quantities are real; \(g,f_{\mathcal C}^{13}\), and
\(f_{\mathcal C}^{33}\) are homogeneous of degree four.

For a set \(\mathcal I\) of ordered pairs \((\sigma,P')\), where
\(P'\in\mathcal S_1\cup\mathcal S_3\) is a Pauli label, put
\begin{equation}
 J(\mathcal I)=
 \sum_{(\sigma,P')\in\mathcal I}
 \Re\!\left[\overline{t_{P'}}\,
        \tau\!\left(P'\sigma R_\sigma(T)\right)\right].
 \label{o2gc:J-contraction}
\end{equation}
Here \(t_{P'}\) is the coefficient of \(P'\) in \(T\); \(P'\) is not
the code projector \(P\). The five choices needed below are defined
by the following table. In each row, sum over every active site \(i\)
and every indicated Pauli pair; \(P'_i\) is the factor at that site.
\begin{center}
\begin{tabular}{@{}lll@{}}
\toprule
Sum & \(\sigma\) and \(P'\) & Condition at site \(i\)\\
\midrule
\(J_{3,L}\) &
\(\sigma\in\{X_i,Y_i\},\ P'\in\mathcal S_3\) & \(P'_i=Z\)\\
\(J_{3,Z}\) &
\(\sigma=Z_i,\ P'\in\mathcal S_3\) & \(P'_i\in\{X,Y\}\)\\
\(J_{3,\perp}\) &
\(\sigma\in\{X_i,Y_i\},\ P'\in\mathcal S_3\) &
\(P'_i\in\{X,Y\},\ P'_i\ne\sigma_i\)\\
\(J_{1,\perp}\) &
\(\sigma\in\{X_i,Y_i\},\ P'\in\mathcal S_1\) &
\(P'_i\in\{X,Y\},\ P'_i\ne\sigma_i\)\\
\(J_{1,Z}\) &
\(\sigma=Z_i,\ P'\in\mathcal S_1\) & \(P'_i\in\{X,Y\}\)\\
\bottomrule
\end{tabular}
\end{center}
All five sums vanish if the common-antiunitary defects vanish.

The certificate polynomial is
\begin{align}
 \mathcal Q(T)={}&-\sum_{E\in\mathcal V_3}c_E^2
 +\frac5{12}J_{3,L}+\frac56J_{3,Z}
 +\frac7{12}J_{3,\perp}+\frac12g\nonumber\\
 &+\frac32J_{1,\perp}-\frac32J_{1,Z}
 +\frac{17}{6}f_{\mathcal F_4}^{13}
 +\frac{19}{6}f_{\mathcal F_4}^{33}
 +\frac{25}{6}f_{\mathcal D_3}^{13}
 +\frac13f_{\mathcal D_2}^{33}.
 \label{o2gc:Q}
\end{align}

\begin{lemma}[Exact coefficient matrix and positivity]
\label{o2gc:lem-matrix}
Order the 32 Pauli labels in \(\mathcal S_1\cup\mathcal S_3\)
lexicographically using \(I<X<Y<Z\). Let \(z\) contain all
\(528\) monomials \(t_{P'}t_{Q'}\), \(P'\le Q'\), with no
\(\sqrt2\) factors on unequal pairs. There is an integer real
symmetric \(528\times528\) matrix \(M\) such that, identically for
every complex \(T\) in this coordinate space,
\begin{equation}
 \mathcal Q(T)=\frac1{12}z^\dagger Mz.
 \label{o2gc:matrix-identity}
\end{equation}
The \(36\) rows and columns indexed by two \(\mathcal S_1\)
coordinates are zero. On the remaining \(492\) coordinates,
\(M\) has \(111\) blocks and satisfies
\begin{equation}
 M_{\mathrm{remaining}}-4I_{492}>0.
 \label{o2gc:matrix-positive}
\end{equation}
In particular,
\begin{equation}
 \mathcal Q(T)\ge\frac13
 \sum_{\substack{P'\le Q'\\
 P'\in\mathcal S_3\ \text{or}\ Q'\in\mathcal S_3}}
 |t_{P'}t_{Q'}|^2.
 \label{o2gc:coercive}
\end{equation}
\end{lemma}

\begin{proof}[Exact finite verification]
Treat the \(t_{P'}\) and their formal conjugates as independent
variables and expand \(12\mathcal Q\) in bidegree \((2,2)\).
Each monomial has the unique form
\(\overline{t_{P'}t_{Q'}}t_{R'}t_{S'}\), with \(P'\le Q'\) and
\(R'\le S'\). Accumulating all ordered trace contributions into
these unordered pair indices uniquely determines the coefficient
matrix \(M\); it is not a fitted Gram representation.

The exact verifier multiplies Paulis with phases \(1,-1,i,-i\),
expands the normalized traces and real parts in
\eqref{o2gc:scalar-contractions}--\eqref{o2gc:Q}, and checks integral
entries and \(M=M^T\). In particular it proves
\eqref{o2gc:matrix-identity} before imposing any encoder equations.
No floating-point eigenvalues or numerical matrix entries enter
this verification.

Conjugation by a four-qubit Pauli multiplies each \(t_{P'}\) by its
Pauli-conjugation sign and leaves \(\mathcal Q\) invariant. Hence a
matrix entry between the monomials \(t_{P'}t_{Q'}\) and
\(t_{R'}t_{S'}\) can be nonzero only when the Pauli labels of
\(P'Q'\) and \(R'S'\) agree modulo phase. This gives the character
blocks. The verifier checks every off-block zero entry, all \(36\)
zero rows and columns, and every remaining block.

The exact characteristic polynomials of all blocks are summarized
in Table~\ref{o2gc:tab-blocks}. Every linear root is greater than
four. Each quadratic factor \(x^2-bx+c\) has real roots because it
is a factor of the characteristic polynomial of a real symmetric
matrix. Its shifted roots \(y=x-4\) have positive sum \(b-8\) and
positive product \(c-4b+16\), as is checked from every listed factor.
Thus both roots exceed four. This proves
\eqref{o2gc:matrix-positive} once the exact block polynomials have
been verified.

The standard-library verifier also establishes positivity without
factoring characteristic polynomials. For each of all \(111\)
integer blocks, it subtracts \(4I\) and performs exact rational
\(LDL^T\) elimination. Every pivot is strictly positive. Explicitly,
at pivot \(d>0\), the remaining Schur complement is updated by
\[
 A_{ij}\longmapsto A_{ij}-\frac{A_{ik}A_{kj}}{d}.
\]
All entries and pivots are rational, and positive pivots give positive
definiteness by congruence. Every block is checked individually, so
this argument does not assume that one representative block covers
its entire symmetry class.

Finally, the zero coordinates are precisely the
\(\binom{8+1}{2}=36\) products of two weight-one coordinates.
On the remaining coordinates \eqref{o2gc:matrix-positive} implies
\eqref{o2gc:coercive}. The matrix is defined completely by
\eqref{o2gc:Q} and the specified monomial ordering; the accompanying
JSON additionally records every block and every positive pivot.
\end{proof}

\begin{table}[htbp]
\centering
\small
\renewcommand{\arraystretch}{1.2}
\begin{tabular}{@{}lrrp{0.64\linewidth}@{}}
\toprule
Character pattern & Size & Count & Characteristic polynomial of a block of \(M\)\\
\midrule
\(I\)&24&1&
\((x-424)(x-192)^4(x-24)^{11}(x-16)^2
(x^2-48x+192)^3\)\\
\(Z\)&3&4&\((x-384)(x-48)^2\)\\
\(pq\)&3&24&\((x-52)(x^2-460x+47872)\)\\
\(pqZ\)&4&48&\((x-50)(x-46)(x^2-248x+10716)\)\\
\(ZZ\)&8&6&
\((x-32)(x-24)(x^2-408x+9856)(x^2-352x+9216)^2\)\\
\(ZZZ\)&12&4&
\((x-136)(x-48)(x^2-352x+9216)^3(x^2-64x+720)^2\)\\
\(pqZZ\)&4&24&\((x-56)(x-44)(x^2-380x+23936)\)\\
\bottomrule
\end{tabular}
\caption{All exact character blocks after removal of the \(36\) zero
\(\mathcal S_1\mathcal S_1\) coordinates. Here \(p,q\in\{X,Y\}\);
all distinct-site and transverse-axis choices are included.
The counts sum to \(111\), and the sizes with multiplicity sum to
\(492\). Each displayed polynomial is verified for every block in
the corresponding class.}
\label{o2gc:tab-blocks}
\end{table}

\begin{proof}[Proof of Proposition~\ref{o2gc:prop-rigidity}]
Under \eqref{o2gc:rigidity-hypotheses}, \(g\) and all five \(J\)-sums
vanish. Each forbidden-set contraction also vanishes because its
factor \(c_E\) is zero. Hence \eqref{o2gc:Q} and
\eqref{o2gc:coercive} give
\[
 -\sum_{E\in\mathcal V_3}c_E^2
 =\mathcal Q(T)
 \ge\frac13
 \sum_{\substack{P'\le Q'\\
 P'\in\mathcal S_3\ \text{or}\ Q'\in\mathcal S_3}}
 |t_{P'}t_{Q'}|^2\ge0.
\]
The left side is nonpositive, so every \(c_E\),
\(E\in\mathcal V_3\), vanishes. The right side includes
\(|t_{P'}|^4\) for every \(P'\in\mathcal S_3\), so all those
coordinates vanish as well. Thus \(T_3=0\).
\end{proof}

\subsection{Completion of the converse and exact attainment}
\label{o2gc:subsec-finish}

\begin{proof}[Proof of Theorem~\ref{o2gc:thm-image}]
For every exact orbit-2 encoder, the analytic construction gives
\eqref{o2gc:T-support}, the partial-isometry identities, and all twelve
physical equations \eqref{o2gc:physical-J}. As noted after
Proposition~\ref{o2gc:prop-rigidity}, these imply every hypothesis
\eqref{o2gc:rigidity-hypotheses}. The proposition therefore gives
\(c_E=0\) for \(E\in\mathcal V_3\). By
\eqref{o2gc:normalization}, the coefficient of \(E\) in
\(K=2TT^\dagger-I\) is \(2c_E\), so
\begin{equation}
 A_3=\frac1{16}\operatorname{Tr}(K_3^2)
     =4\sum_{E\in\mathcal V_3}c_E^2=0.
 \label{o2gc:A3-zero}
\end{equation}
Equation~\eqref{o2gc:plane} gives \(\mathcal L(\alpha)\), with
\(0\le\alpha\le1\) by \(A_1,A_2\ge0\).

For completeness, the reverse inclusion has the following explicit
split-qubit construction. Start with the standard perfect five-qubit
code with stabilizer generators
\[
 XZZXI,\quad IXZZX,\quad XIXZZ,\quad ZXIXZ.
\]
It is pure: every nonidentity Pauli of weight at most two has zero
compression. Its logical \(Z^{\otimes5}\), multiplied by the first
generator, gives the weight-three logical representative \(-YIIYZ\).
Local Clifford rotations and a permutation therefore give a pure
five-qubit code with a chosen logical \(Z_1Z_2Z_3\).

For \(0\le\theta\le\pi/2\), put \(a=\cos(\theta/2)\),
\(b=\sin(\theta/2)\), and split original site \(1\) by the isometry
\[
 W_\theta\ket0=a\ket{00}+b\ket{11},\qquad
 W_\theta\ket1=a\ket{01}+b\ket{10}.
\]
Its exact identities are
\begin{equation}
 W_\theta W_\theta^\dagger
 =\frac12(I+\cos\theta\,Z\otimes I+\sin\theta\,X\otimes X),
 \qquad (Z\otimes Z)W_\theta=W_\theta Z.
 \label{o2gc:split-construction}
\end{equation}
Take the two split outputs as \(A_1,A_2\), original sites \(2,3\)
as \(A_3,A_4\), and original sites \(4,5\) as \(B_1,B_2\).
The original weight-three logical lifts to \(I_{\mathsf B}D\),
which is precisely \eqref{o2gc:record}.

Compression of any six-qubit error of weight at most two through
\(W_\theta\otimes I^{\otimes4}\) is supported on at most two original
sites. Purity of the five-qubit code therefore proves all KL equations.
If the error has a nonidentity factor on an unsplit site, its scalar
compression is zero. On the split pair its scalar is the identity
coefficient \(\operatorname{Tr}(W_\theta^\dagger E W_\theta)/2\).
Consequently the only nonzero weight-one and weight-two coefficients
are \(\cos\theta\) on \(Z_{A_1}\) and \(\sin\theta\) on
\(X_{A_1}X_{A_2}\). Thus
\[
 A_1=\cos^2\theta=\alpha,\qquad A_2=\sin^2\theta=1-\alpha.
\]
The universal inequality \eqref{o2gc:shadow} gives \(A_3=0\), and
\eqref{o2gc:MacWilliams} gives \(\mathcal L(\alpha)\). Every
\(\alpha\in[0,1]\), including both endpoints, occurs. The quantum
Singleton bound excludes distance at least four at length six and
code dimension two, so these are exact \(((6,2,3))\) realizations.
\end{proof}

\begin{remark}[Scope of the converse]
The proof uses no division by a computational coefficient, row norm,
polarization, Schmidt eigenvalue, or Gram minor. It covers arbitrary
interactions among different \(Z\)-pairs in \(T_3\) and both regular
and singular frame configurations. The necessary common-antiunitary
equations already suffice for this enumerator conclusion; no
sufficiency claim for the matrix completion problem is required.
The theorem determines \(\mathcal W_2\), not all code projectors up to
local unitaries, and does not assert global connectedness or uniqueness
of the split-code family. It also does not prove a stronger
partial-isometry inequality with the common-antiunitary equations
omitted.
\end{remark}

\subsection{Exact certificate files and reproducibility}
\label{o2gc:subsec-reproducibility}

The proof and its exact certificate are distributed together in
\path{orbit2_enumerator_package.zip}. All filenames below are relative
to the root of the unpacked package:
\begin{itemize}
\item \path{orbit2_enumerator_section.tex}:
this include-ready section, using the supplement's existing preamble;
\item \path{orbit2_enumerator_section_preview.tex}:
a standalone document containing the same section and its full preamble;
\item \path{orbit2_exact_certificate.py}:
the standard-library verifier that regenerates the integer matrix
from the quartic polynomial and checks all rational pivots;
\item \path{orbit2_certificate_blocks.json}:
all integer blocks, monomial indices, and positive rational pivots;
\item \path{orbit2_verification_log.txt}:
the successful exact-verification record regenerated for this package;
\item \path{README.txt}:
instructions for compiling, including, and verifying the proof;
\item \path{SHA256SUMS}:
SHA-256 file hashes for both LaTeX sources, the verifier, the certificate
JSON, the verification log, and the README.
\end{itemize}
Run from the unpacked package directory with Python~3, without
third-party packages:
\begin{verbatim}
python3 orbit2_exact_certificate.py
python3 orbit2_exact_certificate.py --json regenerated_blocks.json
\end{verbatim}
The regenerated JSON agrees with the supplied
\path{orbit2_certificate_blocks.json}. The verifier reconstructs the
matrix directly from the explicitly coded polynomial identity; it does
not read the LaTeX files or trust a precomputed positive-semidefinite
matrix. The JSON and log are supporting records rather than runtime
dependencies.
The canonical sparse integer matrix has \(2{,}844\) nonzero entries.
Its SHA-256 under the verifier's specified sparse-entry serialization is
\begin{center}
\small\ttfamily
64b192ccce90edf4b583285307f46265435a51b032f72e6493057138a1820b1d
\end{center}
Discovery of the multipliers can use numerical optimization; the proof
does not. Its finite verification uses exact integer coefficient
expansion and rational \texttt{Fraction} arithmetic. The analytic
encoder reduction is proved above, so this section does not depend
on an unproved converse or on a regular-support assumption from an
earlier problem note.

% This file has no document class, preamble, or document environment.
% It uses only the supplement's existing packages, macros, and theorem styles.
% Every new label has the prefix o2gc:.
% Existing abstract/status/count statements elsewhere need separate updating;
% this fragment does not edit or silently replace those statements.

\section{Exact weight-enumerator image of orbit 2}
\label{o2gc:sec}

Orbit~2 is the primitive order-two record
\[
 (q,m,a)=\bigl(2,2,(0,0,1,1,1,1)\bigr).
\]
For the proof, relabel the physical sites in the order
\(B_1,B_2,A_1,A_2,A_3,A_4\). Write
\(\mathsf B=\{B_1,B_2\}\) for the inactive pair and
\(\mathsf A=\{A_1,A_2,A_3,A_4\}\) for the active sites. Its chosen physical
logical involution is
\begin{equation}
 D_{\mathrm{phys}}=D(a;2)=I_{\mathsf B}\otimes D,
 \qquad D=Z_{A_1}Z_{A_2}Z_{A_3}Z_{A_4}.
 \label{o2gc:record}
\end{equation}
Let \(V:\C^2\to(\C^2)^{\otimes6}\) be an isometric encoder, with
code projector \(P=VV^\dagger\), satisfying
\begin{equation}
 D_{\mathrm{phys}}V=VZ_L,\qquad
 V^\dagger EV=\lambda_EI_2\quad(\wt(E)\le2),
 \label{o2gc:KL}
\end{equation}
where \(E\) ranges over Hermitian Pauli strings. For all Pauli weights
put \(\lambda_E=\operatorname{Tr}(PE)/2\); this agrees with the KL scalar
at weights at most two but does not assert scalar compression at higher
weights. We use the supplement's normalization
\[
 A_w=\frac14\sum_{\wt(E)=w}|\operatorname{Tr}(PE)|^2
     =\sum_{\wt(E)=w}\lambda_E^2.
\]

\begin{theorem}[Complete orbit-2 enumerator image]
\label{o2gc:thm-image}
Every exact orbit-2 realization has \(A_3=0\). Its complete
Shor--Laflamme \(A\)-enumerator image is
\begin{equation}
 \boxed{
 \mathcal W_2=
 \left\{\mathcal L(\alpha):0\le\alpha\le1\right\},
 \qquad
 \mathcal L(\alpha)=
 (1,\alpha,1-\alpha,0,11+4\alpha,16-\alpha,3-3\alpha).
 }
 \label{o2gc:image}
\end{equation}
The converse includes every zero-support and singular stratum. No
stabilizer assumption is made.
\end{theorem}

The proof has an analytic reduction and an exact finite certificate.
First, KL and complementary purities force a flat inactive pair and a
specific rank-eight active projector. Mixed KL supplies one common
antiunitary, whose spectral saturation restricts an associated partial
isometry to 32 complex Pauli coordinates. A homogeneous quartic
certificate then forces its weight-three coordinates, and the code's
entire \(A_3\), to vanish. The certificate identity and positivity are
verified with integer and rational arithmetic.

\subsection{Flatness of the inactive pair and the active reflection}
\label{o2gc:subsec-flatness}

Put \(\rho=P/2\). The universal six-qubit rank-two MacWilliams normal
form, with \(x=A_1,y=A_2,z=A_3\), is
\begin{equation}
 (A_0,\ldots,A_6)=
 (1,x,y,z,15-4y-3z,12+3x+4y+3z,4-4x-y-z).
 \label{o2gc:MacWilliams}
\end{equation}
It follows from the quantum MacWilliams transform with factor \(1/32\)
and \(A_j=B_j\) for \(j=0,1,2\). With full six-qubit spin flip
\(\widetilde P=Y^{\otimes6}P^TY^{\otimes6}\), positivity gives
\begin{equation}
 0\le\operatorname{Tr}(P\widetilde P)
 =\frac1{16}\sum_{w=0}^6(-1)^wA_w
 =\frac12(1-A_1-A_2-A_3).
 \label{o2gc:shadow}
\end{equation}

\begin{lemma}[Universal flatness and weight-three support]
\label{o2gc:lem-flatness}
Every encoder in \eqref{o2gc:KL} satisfies
\begin{equation}
 \rho_{\mathsf B}=\frac{I_4}{4},\qquad
 \rho_{\mathsf B A_i}=\frac{I_4}{4}\otimes\rho_{A_i}
 \quad(1\le i\le4),\qquad A_1+A_2+A_3=1.
 \label{o2gc:flatness}
\end{equation}
The only possibly nonzero weight-three coefficients are
\(\lambda_{p_iq_jZ_k}\) on distinct active sites, with
\(p,q\in\{X,Y\}\). Moreover,
\begin{equation}
 \rho_{\mathsf A}=\frac{\Pi}{8},\qquad
 \Pi^\dagger=\Pi=\Pi^2,\qquad \rank\Pi=8,\qquad [\Pi,D]=0,
 \label{o2gc:active-projector}
\end{equation}
and each active parity block of \(\Pi\) has rank four.
\end{lemma}

\begin{proof}
Purify the maximally mixed code with a reference qubit \(R\):
\[
 \ket{\Psi}=
 \frac{\ket0_R\ket{0_L}+\ket1_R\ket{1_L}}{\sqrt2}.
\]
The operator \(G=Z_RI_{\mathsf B}D\) fixes this state. KL says that every
correlation involving a nonidentity reference Pauli and a physical
Pauli of weight at most two vanishes. Multiplication by \(G\) preserves
squared correlations.

Physical coefficients with an odd number of active \(X/Y\) factors
vanish by \(D\)-parity. An active-only coefficient with at least two
\(Z\) factors also vanishes: multiplying by \(G\) gives a
reference-\(Z\) correlation with physical weight \(4-z\le2\), where
\(z\) is the original active \(Z\)-count.
For this purity computation, denote the remaining nonidentity
active-only squared masses by
\[
 \begin{array}{c|cccc}
 \text{mass}&a&u&v&w\\ \hline
 \text{pattern}&Z_i&p_iq_j&p_iq_jZ_k&p_1p_2p_3p_4
 \end{array}
 \qquad(p,q,p_i\in\{X,Y\}).
\]
All displayed sites in a pattern are distinct. Let \(s=a+u+v+w\), and
define the nonnegative masses
\[
 b=\sum_{\substack{F\in\{I,X,Y,Z\}^{\otimes\mathsf B}\\F\ne I}}
       \lambda_F^2,\qquad
 m=\sum_{i=1}^4\sum_{F\ne I}\lambda_{FZ_i}^2.
\]
Here \(F\) in the second sum ranges over all nonidentity spectator
Paulis, including those of spectator weight two.

Purity of a marginal on \(r\) qubits is \(2^{-r}\) times the sum of its
squared Pauli correlations. KL gives
\(\rho_{R\mathsf B}=I_R/2\otimes\rho_{\mathsf B}\).
Equality of complementary purities for \(\mathsf A\) and \(R\mathsf B\)
therefore gives
\begin{equation}
 1+s=2(1+b).
 \label{o2gc:purity-1}
\end{equation}
Let \(h\) be the squared correlation mass on \(R\mathsf A\) with
reference factor \(X\) or \(Y\). Multiplication by \(G\) identifies the
reference-\(Z\) mass with the physical active mass \(1+s\). Comparing
the complementary purities of \(\mathsf B\) and \(R\mathsf A\) gives
\[
 2(1+s)+h=8(1+b),\qquad h=4(1+b).
\]
A surviving reference-\(X/Y\) correlation has odd active transverse
count. If that count is one and its active \(Z\)-count is \(z\), the
term and its \(G\)-multiple have physical weights \(1+z\) and \(4-z\).
At least one weight is at most two, so both correlations vanish.
Thus the surviving transverse count is three. Multiplication by \(G\)
interchanges its physical-weight-three and physical-weight-four
terms, giving
\begin{equation}
 h_3=h_4=\frac h2=2(1+b).
 \label{o2gc:purity-h}
\end{equation}

For each \(i\), compare the complementary marginals
\(\mathsf B A_i\) and \(R(\mathsf A\setminus\{A_i\})\), then sum over
all four active sites. The first marginal's mass is
\[
 1+b+\lambda_{Z_i}^2+\sum_{F\ne I}\lambda_{FZ_i}^2.
\]
On the complementary side, reference identity and reference \(Z\)
together count active strings whose \(i\)th factor is \(I\) or \(Z\).
A string with transverse count \(t\) is therefore counted \(4-t\)
times. Reference-\(X/Y\) physical-weight-three terms appear once and
weight-four terms never appear. Accounting also for the dimensions of
the two marginals gives
\begin{equation}
 2\{4(1+b)+a+m\}
 =4(1+s)-2u-2v-4w+h_3.
 \label{o2gc:purity-2}
\end{equation}
Substitution of \eqref{o2gc:purity-1} and \eqref{o2gc:purity-h} yields
\begin{equation}
 b+m+w=0.
 \label{o2gc:purity-sos}
\end{equation}
Every summand is a sum of real squares. Hence \(b=m=w=0\), every
coefficient they contain vanishes, and \(a+u+v=1\). In particular
\(\rho_{\mathsf B}=I_4/4\). Parity kills spectator-dressed active
\(X/Y\) coefficients, and \(m=0\) kills all spectator-dressed \(Z_i\)
coefficients. This proves the product marginal in
\eqref{o2gc:flatness}.

The full weight-one and weight-two masses are now \(A_1=a\) and
\(A_2=u\). Write \(A_3=v+n\), where \(n\ge0\) collects the remaining
terms supported on two active sites and one spectator site. All other
additional low-weight patterns have vanished. Equation
\eqref{o2gc:shadow} gives \(1+n\le1\), so \(n=0\), \(A_3=v\), and
\(A_1+A_2+A_3=1\).

Finally, KL for all errors on \(\mathsf B\) gives the same marginal
\(I_4/4\) for each logical word. Each word therefore has four equal
Schmidt coefficients \(1/2\). Its active Schmidt space lies in the
corresponding \(D\)-parity sector. Each logical word's active marginal
is its rank-four projector divided by four. Averaging these two
marginals gives \(\Pi/8\), proving
\eqref{o2gc:active-projector}.
\end{proof}

Write the resulting flat parity frames as \(16\times4\) matrices
\(W_0,W_1\):
\begin{equation}
 \ket{z_L}=\frac12\sum_{b=0}^3\ket b_{\mathsf B}\otimes W_z\ket b,
 \qquad W_z^\dagger W_z=I_4,\qquad DW_z=(-1)^zW_z.
 \label{o2gc:frames}
\end{equation}
Their ranges are orthogonal, and
\[
 \Pi=W_0W_0^\dagger+W_1W_1^\dagger
     =4\operatorname{Tr}_{\mathsf B}P.
\]
Set \(K=2\Pi-I_{16}\). Then \(K^\dagger=K\), \(K^2=I_{16}\), and
\([K,D]=0\). For every nonidentity active Pauli \(E\),
\begin{equation}
 \lambda_{I_{\mathsf B}E}
 =\frac18\operatorname{Tr}(\Pi E)
 =\frac1{16}\operatorname{Tr}(KE).
 \label{o2gc:normalization}
\end{equation}
The active support proved above thus gives
\begin{align}
 K&=K_1+K_2+K_3,\label{o2gc:K-support}\\
 K_1&=\sum_{i=1}^4a_iZ_i,\nonumber\\
 K_2&=\sum_{i<j}\sum_{p,q\in\{X,Y\}}u_{ij}^{pq}p_iq_j,\nonumber\\
 K_3&=\sum_{i<j}\sum_{k\notin\{i,j\}}
              \sum_{p,q\in\{X,Y\}}v_{ij;k}^{pq}p_iq_jZ_k.\nonumber
\end{align}
The three sets have \(4,24,48\) real coefficients, respectively. Their
squared masses are
\begin{equation}
 \alpha=\frac1{16}\operatorname{Tr}(K_1^2)=A_1,\qquad
 \beta=\frac1{16}\operatorname{Tr}(K_2^2)=A_2,\qquad
 \nu=\frac1{16}\operatorname{Tr}(K_3^2)=A_3,
 \quad \alpha+\beta+\nu=1.
 \label{o2gc:masses}
\end{equation}
Consequently \eqref{o2gc:MacWilliams} already reduces the enumerator to
\begin{equation}
 (A_0,\ldots,A_6)=
 (1,\alpha,1-\alpha-\nu,\nu,
   11+4\alpha+\nu,16-\alpha-\nu,3-3\alpha).
 \label{o2gc:plane}
\end{equation}
The remaining converse is \(\nu=0\). An arbitrary reflection having
\eqref{o2gc:K-support} need not be completable to a six-qubit code; the
mixed KL constraints are essential in the next step.

\subsection{One common antiunitary from mixed KL}
\label{o2gc:subsec-commonJ}

On the spectator pair choose the fixed basis of
\(\Theta_{\mathsf B}=(Y\otimes Y)\mathcal K\), where \(\mathcal K\)
denotes computational complex conjugation. Explicitly, let \(M_{\mathsf B}\)
have columns
\[
 i\Phi_+,\quad\Phi_-,\quad\Psi_+,\quad i\Psi_-,
 \qquad
 \Phi_\pm=\frac{\ket{00}\pm\ket{11}}{\sqrt2},\quad
 \Psi_\pm=\frac{\ket{01}\pm\ket{10}}{\sqrt2}.
\]
Each column is fixed by \(\Theta_{\mathsf B}\). This is a common
spectator magic basis, not an independently chosen basis for each
observable. In that basis put
\[
 Q_z=W_z\overline{M_{\mathsf B}},\qquad Q=(Q_0\ Q_1).
\]
The conjugate follows from the bipartite expansion
\eqref{o2gc:frames}. We have \(Q^\dagger Q=I_8\) and \(QQ^\dagger=\Pi\).

\begin{lemma}[Common antiunitary equations]
\label{o2gc:lem-commonJ}
For every real single-body active Hamiltonian \(H\), put
\[
 h=Q^\dagger HQ,\qquad
 \mu_H=\frac18\operatorname{Tr}(\Pi H).
\]
Then the same skew-unitary matrix
\begin{equation}
 J_0=\begin{pmatrix}0&I_4\\-I_4&0\end{pmatrix}
 \label{o2gc:J0}
\end{equation}
satisfies
\begin{equation}
 hJ_0+J_0h^T=2\mu_HJ_0.
 \label{o2gc:commonJ}
\end{equation}
\end{lemma}

\begin{proof}
Each single-spectator Pauli \(F\) changes sign under
\(\Theta_{\mathsf B}\). Hence \(M_{\mathsf B}^\dagger F M_{\mathsf B}\)
is imaginary, antisymmetric, and Hermitian. The six such Paulis span
the six-dimensional space of imaginary antisymmetric Hermitian
\(4\times4\) matrices.

For every single-active Pauli \(E\), mixed KL gives, for all six \(F\)
and all \(z,z'\),
\[
 \frac14\operatorname{Tr}\!\left[
 (M_{\mathsf B}^\dagger F M_{\mathsf B})^T
 Q_z^\dagger E Q_{z'}\right]=0.
\]
Orthogonality to this antisymmetric space makes every block
\(Q_z^\dagger EQ_{z'}\) complex symmetric. A diagonal block is also
Hermitian and is therefore real symmetric.

The stronger product-marginal identity in \eqref{o2gc:flatness}, not
merely spectator flatness, gives
\[
 h_{00}+h_{11}=2\mu_HI_4.
\]
Indeed, tracing \((I_{\mathsf B}\otimes H)\rho\) over the active sites
gives the spectator matrix \((h_{00}+h_{11})^T/8\) in this basis;
the product marginal instead gives \(\mu_HI_4/4\).
Thus
\[
 h=\begin{pmatrix}
 R&C\\ \overline C&2\mu_HI_4-R
 \end{pmatrix},
 \qquad R=R^T=\overline R,\qquad C=C^T.
\]
Direct block multiplication proves \eqref{o2gc:commonJ}.
\end{proof}

The antiunitary \(J_0\mathcal K\) squares to \(-I_8\) and reverses
every centered compression \(h-\mu_HI_8\). The same \(J_0\) works for
all twelve single-active Paulis. Separate spectral pairings, or a
condition using only the four \(Z_i\), would not provide the full
hypotheses used in the final certificate.

\subsection{A sparse partial isometry on four active qubits}
\label{o2gc:subsec-T}

Define
\begin{equation}
 \Omega=QJ_0Q^T=Q_0Q_1^T-Q_1Q_0^T,\qquad
 S=Y^{\otimes4},\qquad T=\Omega S.
 \label{o2gc:T-definition}
\end{equation}
Here \(S\) is real, symmetric, and unitary. Using
\(Q^T\overline Q=I_8\) gives
\[
 \Omega^T=-\Omega,\qquad
 \Omega\Omega^\dagger=\Pi,\qquad
 \Omega^\dagger\Omega=\Pi^T.
\]
Therefore
\begin{equation}
 TT^\dagger=\Pi,\qquad
 T^\dagger T=\widetilde\Pi:=S\Pi^TS,\qquad
 ST^TS=-T,\qquad \operatorname{Tr}(TT^\dagger)=8.
 \label{o2gc:partial-isometry}
\end{equation}
The Pauli identity \(SE^TS=(-1)^{\wt(E)}E\) means that only odd
Pauli weights can occur in \(T\).

\begin{lemma}[Exact edge sparsity]
\label{o2gc:lem-edges}
For computational active strings \(x,y\in\{0,1\}^4\),
\[
 d_H(x,y)=1\quad\Longrightarrow\quad\Omega_{xy}=0.
\]
In fact \(\Omega\) is supported only at Hamming distance three, and
\(T\) only at Hamming distance one.
\end{lemma}

\begin{proof}
Parity already makes \(\Omega\) zero between equal-parity strings,
since it pairs \(Q_0\) and \(Q_1\). Fix an even string \(x\), set
\(s_i=(-1)^{x_i}\), \(t=\sum_i a_is_i\), and write \(\bar x\) for
its complement. Consider
\[
 H_x=\sum_i s_iZ_i,\qquad
 r_x=\bra xQ_0,\qquad r_{\bar x}=\bra{\bar x}Q_0.
\]
The allowed reflection support \eqref{o2gc:K-support} gives
\[
 \|r_x\|^2=\Pi_{xx}=\frac{1+t}{2},\qquad
 \|r_{\bar x}\|^2=\frac{1-t}{2},\qquad
 r_xr_{\bar x}^\dagger=\Pi_{x\bar x}=0.
\]
The last equality uses the absence of any term flipping all four
active bits.

The eigenvalue of \(H_x\) on \(y\) is \(4-2d_H(x,y)\). Thus its
even-parity restriction is
\(4(\ket x\bra x-\ket{\bar x}\bra{\bar x})\), and its odd-parity
spectrum is \(\{-2,+2\}\). With \(h_z=Q_z^\dagger H_xQ_z\), this gives
\[
 h_0=4r_x^\dagger r_x-4r_{\bar x}^\dagger r_{\bar x},\qquad
 h_0^Tr_x^T=2(1+t)r_x^T.
\]
Moreover \(\mu_{H_x}=t\). Equation~\eqref{o2gc:commonJ} implies
\[
 h_1=2tI_4-h_0^T,\qquad h_1r_x^T=-2r_x^T.
\]
Set \(w=Q_1r_x^T\). Then
\[
 \bra w(H_x+2I)\ket w=0.
\]
On odd parity \(H_x+2I\) is positive semidefinite, with kernel
consisting exactly of the strings at distance three from \(x\).
Consequently \(w\) has no distance-one components. For even \(x\)
and odd \(y\),
\[
 \Omega_{xy}=r_x(\bra yQ_1)^T=\bra y w,
\]
so the desired entries vanish. If \(r_x=0\), that row already vanishes;
no division is needed. Skew symmetry supplies the reverse ordering.
Only odd distances are possible, so the remaining distance is three.
Right multiplication by \(S\) complements the column string up to a
phase, taking distance three to distance one.
\end{proof}

An operator supported only on distance-one matrix entries has exactly
one \(X/Y\) factor in its Pauli expansion. Combining this with the odd
Pauli weights in \eqref{o2gc:partial-isometry} gives
\begin{align}
 T&=T_1+T_3,\label{o2gc:T-support}\\
 T_1&=\sum_i(t_i^XX_i+t_i^YY_i),\nonumber\\
 T_3&=\sum_i\sum_{\substack{j<k\\j,k\ne i}}
          (t_{i;jk}^XX_i+t_{i;jk}^YY_i)Z_jZ_k.\nonumber
\end{align}
There are \(8+24=32\) complex coefficients.

For each single-active Pauli \(\sigma\), multiply
\eqref{o2gc:commonJ} by \(Q,Q^T\) to obtain
\[
 \Pi\sigma\Omega+\Omega\sigma^T\Pi^T=2\mu_\sigma\Omega.
\]
Since \(S\sigma^TS=-\sigma\), multiplication on the right by \(S\)
yields the physical common-antiunitary equation
\begin{equation}
 \Pi\sigma T-T\sigma\widetilde\Pi=2\mu_\sigma T,\qquad
 \mu_\sigma=\frac18\operatorname{Tr}(\Pi\sigma).
 \label{o2gc:physical-J}
\end{equation}
All steps so far are valid for arbitrary computational zeros and
singular overlap patterns.

\subsection{A homogeneous rigidity theorem}
\label{o2gc:subsec-rigidity}

The remaining finite statement can be formulated independently of
encoders. All operators in this subsection act on the four active
qubits; active identity factors are suppressed. Write
\[
 \tau(M)=\frac1{16}\operatorname{Tr}(M),\qquad
 \mathcal S_1=\{X_i,Y_i:1\le i\le4\},
\]
\[
 \mathcal S_3=\{p_iZ_jZ_k:p\in\{X,Y\},\ j<k,\ i\notin\{j,k\}\}.
\]
For an arbitrary \(T\) in the complex span of
\(\mathcal S_1\cup\mathcal S_3\), let \(T_1,T_3\) denote these two
coordinate parts and put
\begin{equation}
 N=\tau(TT^\dagger),\qquad c_E=\tau(ETT^\dagger).
 \label{o2gc:homogeneous-scalars}
\end{equation}
Each \(c_E\) is real. Define the homogeneous cubic defects
\begin{align}
 F(T)&=TT^\dagger T-2NT,\label{o2gc:defects}\\
 R_\sigma(T)&=TT^\dagger\sigma T-T\sigma T^\dagger T-4c_\sigma T.
 \nonumber
\end{align}
The target and forbidden Pauli sets are
\begin{align}
 \mathcal V_3&=\{p_iq_jZ_k:i<j,\ k\notin\{i,j\},\
                                  p,q\in\{X,Y\}\},\nonumber\\
 \mathcal D_2&=\{Z_iZ_j:i<j\},\nonumber\\
 \mathcal D_3&=\{Z_iZ_jZ_k:i<j<k\},\nonumber\\
 \mathcal F_4&=\{p_iq_jZ_kZ_\ell:i<j,\
                \{k,\ell\}=\{1,2,3,4\}\setminus\{i,j\},\
                p,q\in\{X,Y\}\}.
 \label{o2gc:Pauli-sets}
\end{align}
Their sizes are \(48,6,4,24\), respectively.

\begin{proposition}[Global common-antiunitary rigidity]
\label{o2gc:prop-rigidity}
Suppose \(T\) has the support \eqref{o2gc:T-support} and
\begin{equation}
 F(T)=0,\qquad R_\sigma(T)=0
 \quad\text{for all twelve single-active Paulis},\qquad
 c_E=0\quad(E\in\mathcal D_2\cup\mathcal D_3\cup\mathcal F_4).
 \label{o2gc:rigidity-hypotheses}
\end{equation}
Then \(T_3=0\) and \(c_E=0\) for every \(E\in\mathcal V_3\).
\end{proposition}

For the matrix constructed from an orbit-2 code,
\(N=\tau(\Pi)=1/2\), so \(F(T)=0\) follows from its partial-isometry
identities. Also \(c_\sigma=\mu_\sigma/2\), so
\eqref{o2gc:physical-J} gives \(R_\sigma(T)=0\) with exactly the
factor \(4\) in \eqref{o2gc:defects}. Finally, the support of
\(K=2TT^\dagger-I\) gives all the forbidden zeros in
\eqref{o2gc:rigidity-hypotheses}. Thus the proposition applies once
its certificate is proved.

\subsection{The exact positive quartic certificate}
\label{o2gc:subsec-certificate}

Define the scalar defect contractions
\begin{align}
 g&=\Re\,\tau(T_3^\dagger F(T)),\nonumber\\
 c_E^{13}&=\tau\!\left(E(T_1T_3^\dagger+T_3T_1^\dagger)\right),
 \qquad c_E^{33}=\tau(ET_3T_3^\dagger),\nonumber\\
 f_{\mathcal C}^{13}&=\sum_{E\in\mathcal C}c_Ec_E^{13},\qquad
 f_{\mathcal C}^{33}=\sum_{E\in\mathcal C}c_Ec_E^{33}.
 \label{o2gc:scalar-contractions}
\end{align}
The last two definitions apply to any Pauli set \(\mathcal C\).
All displayed quantities are real; \(g,f_{\mathcal C}^{13}\), and
\(f_{\mathcal C}^{33}\) are homogeneous of degree four.

For a set \(\mathcal I\) of ordered pairs \((\sigma,P')\), where
\(P'\in\mathcal S_1\cup\mathcal S_3\) is a Pauli label, put
\begin{equation}
 J(\mathcal I)=
 \sum_{(\sigma,P')\in\mathcal I}
 \Re\!\left[\overline{t_{P'}}\,
        \tau\!\left(P'\sigma R_\sigma(T)\right)\right].
 \label{o2gc:J-contraction}
\end{equation}
Here \(t_{P'}\) is the coefficient of \(P'\) in \(T\); \(P'\) is not
the code projector \(P\). The five choices needed below are defined
by the following table. In each row, sum over every active site \(i\)
and every indicated Pauli pair; \(P'_i\) is the factor at that site.
\begin{center}
\begin{tabular}{@{}lll@{}}
\toprule
Sum & \(\sigma\) and \(P'\) & Condition at site \(i\)\\
\midrule
\(J_{3,L}\) &
\(\sigma\in\{X_i,Y_i\},\ P'\in\mathcal S_3\) & \(P'_i=Z\)\\
\(J_{3,Z}\) &
\(\sigma=Z_i,\ P'\in\mathcal S_3\) & \(P'_i\in\{X,Y\}\)\\
\(J_{3,\perp}\) &
\(\sigma\in\{X_i,Y_i\},\ P'\in\mathcal S_3\) &
\(P'_i\in\{X,Y\},\ P'_i\ne\sigma_i\)\\
\(J_{1,\perp}\) &
\(\sigma\in\{X_i,Y_i\},\ P'\in\mathcal S_1\) &
\(P'_i\in\{X,Y\},\ P'_i\ne\sigma_i\)\\
\(J_{1,Z}\) &
\(\sigma=Z_i,\ P'\in\mathcal S_1\) & \(P'_i\in\{X,Y\}\)\\
\bottomrule
\end{tabular}
\end{center}
All five sums vanish if the common-antiunitary defects vanish.

The certificate polynomial is
\begin{align}
 \mathcal Q(T)={}&-\sum_{E\in\mathcal V_3}c_E^2
 +\frac5{12}J_{3,L}+\frac56J_{3,Z}
 +\frac7{12}J_{3,\perp}+\frac12g\nonumber\\
 &+\frac32J_{1,\perp}-\frac32J_{1,Z}
 +\frac{17}{6}f_{\mathcal F_4}^{13}
 +\frac{19}{6}f_{\mathcal F_4}^{33}
 +\frac{25}{6}f_{\mathcal D_3}^{13}
 +\frac13f_{\mathcal D_2}^{33}.
 \label{o2gc:Q}
\end{align}

\begin{lemma}[Exact coefficient matrix and positivity]
\label{o2gc:lem-matrix}
Order the 32 Pauli labels in \(\mathcal S_1\cup\mathcal S_3\)
lexicographically using \(I<X<Y<Z\). Let \(z\) contain all
\(528\) monomials \(t_{P'}t_{Q'}\), \(P'\le Q'\), with no
\(\sqrt2\) factors on unequal pairs. There is an integer real
symmetric \(528\times528\) matrix \(M\) such that, identically for
every complex \(T\) in this coordinate space,
\begin{equation}
 \mathcal Q(T)=\frac1{12}z^\dagger Mz.
 \label{o2gc:matrix-identity}
\end{equation}
The \(36\) rows and columns indexed by two \(\mathcal S_1\)
coordinates are zero. On the remaining \(492\) coordinates,
\(M\) has \(111\) blocks and satisfies
\begin{equation}
 M_{\mathrm{remaining}}-4I_{492}>0.
 \label{o2gc:matrix-positive}
\end{equation}
In particular,
\begin{equation}
 \mathcal Q(T)\ge\frac13
 \sum_{\substack{P'\le Q'\\
 P'\in\mathcal S_3\ \text{or}\ Q'\in\mathcal S_3}}
 |t_{P'}t_{Q'}|^2.
 \label{o2gc:coercive}
\end{equation}
\end{lemma}

\begin{proof}[Exact finite verification]
Treat the \(t_{P'}\) and their formal conjugates as independent
variables and expand \(12\mathcal Q\) in bidegree \((2,2)\).
Each monomial has the unique form
\(\overline{t_{P'}t_{Q'}}t_{R'}t_{S'}\), with \(P'\le Q'\) and
\(R'\le S'\). Accumulating all ordered trace contributions into
these unordered pair indices uniquely determines the coefficient
matrix \(M\); it is not a fitted Gram representation.

The exact verifier multiplies Paulis with phases \(1,-1,i,-i\),
expands the normalized traces and real parts in
\eqref{o2gc:scalar-contractions}--\eqref{o2gc:Q}, and checks integral
entries and \(M=M^T\). In particular it proves
\eqref{o2gc:matrix-identity} before imposing any encoder equations.
No floating-point eigenvalues or numerical matrix entries enter
this verification.

Conjugation by a four-qubit Pauli multiplies each \(t_{P'}\) by its
Pauli-conjugation sign and leaves \(\mathcal Q\) invariant. Hence a
matrix entry between the monomials \(t_{P'}t_{Q'}\) and
\(t_{R'}t_{S'}\) can be nonzero only when the Pauli labels of
\(P'Q'\) and \(R'S'\) agree modulo phase. This gives the character
blocks. The verifier checks every off-block zero entry, all \(36\)
zero rows and columns, and every remaining block.

The exact characteristic polynomials of all blocks are summarized
in Table~\ref{o2gc:tab-blocks}. Every linear root is greater than
four. Each quadratic factor \(x^2-bx+c\) has real roots because it
is a factor of the characteristic polynomial of a real symmetric
matrix. Its shifted roots \(y=x-4\) have positive sum \(b-8\) and
positive product \(c-4b+16\), as is checked from every listed factor.
Thus both roots exceed four. This proves
\eqref{o2gc:matrix-positive} once the exact block polynomials have
been verified.

The standard-library verifier also establishes positivity without
factoring characteristic polynomials. For each of all \(111\)
integer blocks, it subtracts \(4I\) and performs exact rational
\(LDL^T\) elimination. Every pivot is strictly positive. Explicitly,
at pivot \(d>0\), the remaining Schur complement is updated by
\[
 A_{ij}\longmapsto A_{ij}-\frac{A_{ik}A_{kj}}{d}.
\]
All entries and pivots are rational, and positive pivots give positive
definiteness by congruence. Every block is checked individually, so
this argument does not assume that one representative block covers
its entire symmetry class.

Finally, the zero coordinates are precisely the
\(\binom{8+1}{2}=36\) products of two weight-one coordinates.
On the remaining coordinates \eqref{o2gc:matrix-positive} implies
\eqref{o2gc:coercive}. The matrix is defined completely by
\eqref{o2gc:Q} and the specified monomial ordering; the accompanying
JSON additionally records every block and every positive pivot.
\end{proof}

\begin{table}[htbp]
\centering
\small
\renewcommand{\arraystretch}{1.2}
\begin{tabular}{@{}lrrp{0.64\linewidth}@{}}
\toprule
Character pattern & Size & Count & Characteristic polynomial of a block of \(M\)\\
\midrule
\(I\)&24&1&
\((x-424)(x-192)^4(x-24)^{11}(x-16)^2
(x^2-48x+192)^3\)\\
\(Z\)&3&4&\((x-384)(x-48)^2\)\\
\(pq\)&3&24&\((x-52)(x^2-460x+47872)\)\\
\(pqZ\)&4&48&\((x-50)(x-46)(x^2-248x+10716)\)\\
\(ZZ\)&8&6&
\((x-32)(x-24)(x^2-408x+9856)(x^2-352x+9216)^2\)\\
\(ZZZ\)&12&4&
\((x-136)(x-48)(x^2-352x+9216)^3(x^2-64x+720)^2\)\\
\(pqZZ\)&4&24&\((x-56)(x-44)(x^2-380x+23936)\)\\
\bottomrule
\end{tabular}
\caption{All exact character blocks after removal of the \(36\) zero
\(\mathcal S_1\mathcal S_1\) coordinates. Here \(p,q\in\{X,Y\}\);
all distinct-site and transverse-axis choices are included.
The counts sum to \(111\), and the sizes with multiplicity sum to
\(492\). Each displayed polynomial is verified for every block in
the corresponding class.}
\label{o2gc:tab-blocks}
\end{table}

\begin{proof}[Proof of Proposition~\ref{o2gc:prop-rigidity}]
Under \eqref{o2gc:rigidity-hypotheses}, \(g\) and all five \(J\)-sums
vanish. Each forbidden-set contraction also vanishes because its
factor \(c_E\) is zero. Hence \eqref{o2gc:Q} and
\eqref{o2gc:coercive} give
\[
 -\sum_{E\in\mathcal V_3}c_E^2
 =\mathcal Q(T)
 \ge\frac13
 \sum_{\substack{P'\le Q'\\
 P'\in\mathcal S_3\ \text{or}\ Q'\in\mathcal S_3}}
 |t_{P'}t_{Q'}|^2\ge0.
\]
The left side is nonpositive, so every \(c_E\),
\(E\in\mathcal V_3\), vanishes. The right side includes
\(|t_{P'}|^4\) for every \(P'\in\mathcal S_3\), so all those
coordinates vanish as well. Thus \(T_3=0\).
\end{proof}

\subsection{Completion of the converse and exact attainment}
\label{o2gc:subsec-finish}

\begin{proof}[Proof of Theorem~\ref{o2gc:thm-image}]
For every exact orbit-2 encoder, the analytic construction gives
\eqref{o2gc:T-support}, the partial-isometry identities, and all twelve
physical equations \eqref{o2gc:physical-J}. As noted after
Proposition~\ref{o2gc:prop-rigidity}, these imply every hypothesis
\eqref{o2gc:rigidity-hypotheses}. The proposition therefore gives
\(c_E=0\) for \(E\in\mathcal V_3\). By
\eqref{o2gc:normalization}, the coefficient of \(E\) in
\(K=2TT^\dagger-I\) is \(2c_E\), so
\begin{equation}
 A_3=\frac1{16}\operatorname{Tr}(K_3^2)
     =4\sum_{E\in\mathcal V_3}c_E^2=0.
 \label{o2gc:A3-zero}
\end{equation}
Equation~\eqref{o2gc:plane} gives \(\mathcal L(\alpha)\), with
\(0\le\alpha\le1\) by \(A_1,A_2\ge0\).

For completeness, the reverse inclusion has the following explicit
split-qubit construction. Start with the standard perfect five-qubit
code with stabilizer generators
\[
 XZZXI,\quad IXZZX,\quad XIXZZ,\quad ZXIXZ.
\]
It is pure: every nonidentity Pauli of weight at most two has zero
compression. Its logical \(Z^{\otimes5}\), multiplied by the first
generator, gives the weight-three logical representative \(-YIIYZ\).
Local Clifford rotations and a permutation therefore give a pure
five-qubit code with a chosen logical \(Z_1Z_2Z_3\).

For \(0\le\theta\le\pi/2\), put \(a=\cos(\theta/2)\),
\(b=\sin(\theta/2)\), and split original site \(1\) by the isometry
\[
 W_\theta\ket0=a\ket{00}+b\ket{11},\qquad
 W_\theta\ket1=a\ket{01}+b\ket{10}.
\]
Its exact identities are
\begin{equation}
 W_\theta W_\theta^\dagger
 =\frac12(I+\cos\theta\,Z\otimes I+\sin\theta\,X\otimes X),
 \qquad (Z\otimes Z)W_\theta=W_\theta Z.
 \label{o2gc:split-construction}
\end{equation}
Take the two split outputs as \(A_1,A_2\), original sites \(2,3\)
as \(A_3,A_4\), and original sites \(4,5\) as \(B_1,B_2\).
The original weight-three logical lifts to \(I_{\mathsf B}D\),
which is precisely \eqref{o2gc:record}.

Compression of any six-qubit error of weight at most two through
\(W_\theta\otimes I^{\otimes4}\) is supported on at most two original
sites. Purity of the five-qubit code therefore proves all KL equations.
If the error has a nonidentity factor on an unsplit site, its scalar
compression is zero. On the split pair its scalar is the identity
coefficient \(\operatorname{Tr}(W_\theta^\dagger E W_\theta)/2\).
Consequently the only nonzero weight-one and weight-two coefficients
are \(\cos\theta\) on \(Z_{A_1}\) and \(\sin\theta\) on
\(X_{A_1}X_{A_2}\). Thus
\[
 A_1=\cos^2\theta=\alpha,\qquad A_2=\sin^2\theta=1-\alpha.
\]
The universal inequality \eqref{o2gc:shadow} gives \(A_3=0\), and
\eqref{o2gc:MacWilliams} gives \(\mathcal L(\alpha)\). Every
\(\alpha\in[0,1]\), including both endpoints, occurs. The quantum
Singleton bound excludes distance at least four at length six and
code dimension two, so these are exact \(((6,2,3))\) realizations.
\end{proof}

\begin{remark}[Scope of the converse]
The proof uses no division by a computational coefficient, row norm,
polarization, Schmidt eigenvalue, or Gram minor. It covers arbitrary
interactions among different \(Z\)-pairs in \(T_3\) and both regular
and singular frame configurations. The necessary common-antiunitary
equations already suffice for this enumerator conclusion; no
sufficiency claim for the matrix completion problem is required.
The theorem determines \(\mathcal W_2\), not all code projectors up to
local unitaries, and does not assert global connectedness or uniqueness
of the split-code family. It also does not prove a stronger
partial-isometry inequality with the common-antiunitary equations
omitted.
\end{remark}

\subsection{Exact certificate files and reproducibility}
\label{o2gc:subsec-reproducibility}

The proof and its exact certificate are distributed together in
\path{orbit2_enumerator_package.zip}. All filenames below are relative
to the root of the unpacked package:
\begin{itemize}
\item \path{orbit2_enumerator_section.tex}:
this include-ready section, using the supplement's existing preamble;
\item \path{orbit2_enumerator_section_preview.tex}:
a standalone document containing the same section and its full preamble;
\item \path{orbit2_exact_certificate.py}:
the standard-library verifier that regenerates the integer matrix
from the quartic polynomial and checks all rational pivots;
\item \path{orbit2_certificate_blocks.json}:
all integer blocks, monomial indices, and positive rational pivots;
\item \path{orbit2_verification_log.txt}:
the successful exact-verification record regenerated for this package;
\item \path{README.txt}:
instructions for compiling, including, and verifying the proof;
\item \path{SHA256SUMS}:
SHA-256 file hashes for both LaTeX sources, the verifier, the certificate
JSON, the verification log, and the README.
\end{itemize}
Run from the unpacked package directory with Python~3, without
third-party packages:
\begin{verbatim}
python3 orbit2_exact_certificate.py
python3 orbit2_exact_certificate.py --json regenerated_blocks.json
\end{verbatim}
The regenerated JSON agrees with the supplied
\path{orbit2_certificate_blocks.json}. The verifier reconstructs the
matrix directly from the explicitly coded polynomial identity; it does
not read the LaTeX files or trust a precomputed positive-semidefinite
matrix. The JSON and log are supporting records rather than runtime
dependencies.
The canonical sparse integer matrix has \(2{,}844\) nonzero entries.
Its SHA-256 under the verifier's specified sparse-entry serialization is
\begin{center}
\small\ttfamily
64b192ccce90edf4b583285307f46265435a51b032f72e6493057138a1820b1d
\end{center}
Discovery of the multipliers can use numerical optimization; the proof
does not. Its finite verification uses exact integer coefficient
expansion and rational \texttt{Fraction} arithmetic. The analytic
encoder reduction is proved above, so this section does not depend
on an unproved converse or on a regular-support assumption from an
earlier problem note.
